% !TEX root = ../monograph.tex
\chapter{Insertions and finite-scale comparison}
\label{ch11:insertions}

A step acts linearly on an inserted density even though
its effective-action coordinates evolve nonlinearly. Keeping the density
itself in view resolves two difficulties: several insertions generate
joint terms, and an auxiliary sequence of block kernels can preserve
all original correlation integrals without preserving their intermediate
densities. This chapter establishes that finite-step calculus and the
local comparison estimates it needs. Statements about a limiting quantum
field require further uniform estimates; they are not consequences of
the source algebra alone.

Throughout, the polynomial label space and its local section are
$V_d$ and $P_d^{\rm loc}$ from Chapter~\ref{ch02:observables}.
All block kernels are the normalized kernels of
Chapter~\ref{ch04:blocking}. A \emph{history} means one fixed choice of these
kernels, inverse couplings, adjacent path pairs, and local representations; the local representations are those of Chapter~\ref{ch09:renormalization}, whose activities are averaged over the forests (Section~\ref{ch09:forest-average}). Throughout this chapter we assume Assumption~\ref{ch09:forest-assumption}, which is not proved in this monograph, together with its analogue for the restarted steps of Section~\ref{ch11:restart-construction}, and, for the steps with insertions, its extension stated as Hypothesis~(H) in Section~\ref{ch11:marked-step}, which is not proved either.
The estimates of Chapters~\ref{ch06:large-fields} and
\ref{ch09:renormalization} are used only on that same history.

The argument has four parts. First the squarefree source algebra
and support estimates retain every joint insertion term. Next the
single-insertion estimates and the joint Gaussian comparison construct
the restarted history and the normalized product-moment bounds.
Backward comparison then relates local trials to the
renormalized fields and microscopic equations. Finally a common
localization map onto the polynomial label space and a larger source weight give the bound from a finite separating family, from which uniform finite-difference bounds follow.

\section{An algebra for finitely many insertions}

Let $\mathcal I$ be a finite set of insertion indices and define
\[
 \mathcal A_{\mathcal I}
 =\CC[z_i:i\in\mathcal I]/(z_i^2:i\in\mathcal I),\qquad
 z_S=\prod_{i\in S}z_i.
\]
The variables commute; a monomial with a repeated index is zero.
Every element has the unique form $\sum_{S\subseteq\mathcal I}a_Sz_S$.
If $u$ has zero constant coefficient, then $u^{|\mathcal I|+1}=0$,
so the series for $\exp u$ and $\log(1+u)$ terminate.
For positive weights $\omega_i$, put
\[
 \|a\|_{\omega}=\sum_{S\subseteq\mathcal I}
                         |a_S|\prod_{i\in S}\omega_i.
\]
This norm is submultiplicative by the disjoint-subset convolution
formula. The weights rescale source variables; they are not probabilities.

Write $[z_S]M=M_S$ for the coefficient of $z_S$ in $M$. If
$M_\varnothing=1$ and $S\ne\varnothing$,
\begin{equation}
 [z_S]\log M
 =\sum_{\pi\in\mathcal P(S)}
    (-1)^{|\pi|-1}(|\pi|-1)!\prod_{I\in\pi}M_I .
 \label{ch11:cumulants}
\end{equation}
Here $\mathcal P(S)$ denotes the partitions of $S$ into nonempty blocks,
and $|\pi|$ is the number of blocks of $\pi$.
This is the standard moment--cumulant identity; we use the classical
cumulant calculus of Leonov and Shiryaev~\cite{LeonovShiryaev1959},
rather than develop a second connected-expansion formalism.
Nilpotence makes the formula an algebraic identity even when there
is no positive probability law behind the coefficients.
On bounded coefficient sets, $\log$ is Lipschitz with a constant
depending only on $|\mathcal I|$ and the coefficient bound:
expand the finite sum and telescope each product.

\subsection{The exact density and its integral}

At cutoff $N$, let $\rho_N\ge0$ be the density without insertions, with
$Z_N=\int\rho_N\,d\lambda_N>0$, where $\lambda_N$ is product Haar
probability measure. For dimensionful insertions $F_i$, assume every
$\rho_N\prod_{i\in S}F_i$ is integrable and set
\begin{equation}
 q_N(z,U)=\rho_N(U)\prod_{i\in\mathcal I}(1+z_iF_i(U)),
 \qquad q_m=\mathcal B^{(L_m)}_\beta q_{m+1}.
 \label{ch11:source-density}
\end{equation}
Here $\beta$ is the inverse coupling used for the step $m+1\to m$ on the chosen history. The block map of \eqref{ch04:block-operator} acts coefficientwise. Consequently
$Z_N^{-1}\int[z_S]q_m\,d\lambda_m$ is the original cutoff moment
of the insertions indexed by $S$.

If one instead uses the dimensionless one-insertion density of degree
$d$, its recursion is $H_m=L_m^d\mathcal B^{(L_m)}_\beta H_{m+1}$.
The scalar rescaling belongs to this convention; it is absent from
\eqref{ch11:source-density}, where the dimensionful normalization is
already in $F_i$. Mixing these two conventions changes the identity-field
insertion. In particular the reference-scale coefficient of the dimensionful
identity field in a degree-$d$ convention is $a_{r(d)}^d e_0$, not the coefficient
obtained from the individual scaling dimension of $e_0$.

\begin{proposition}[Changing the block history]
\label{ch11:mass}
Replace any kernels below a fixed scale index $k$ by other normalized positive
block kernels. Transfer the full polynomial $q_k$ through them.
Every coefficient integral is unchanged. A field-independent polynomial
$J\in\mathcal A_{\mathcal I}$ commutes with every such transfer; if
$J_\varnothing=1$, the density without insertions and its partition integral
are unchanged.
\end{proposition}
\begin{proof}
For any integrable coefficient $h$, normalization and Fubini give
$\int\mathcal B h\,d\lambda=\int h\,d\lambda$ and
$\|\mathcal B h\|_1\le\|h\|_1$.
Iteration proves the first assertion. Multiplication by $J$ is a fixed
linear map of the finite coefficient space, so it commutes with the
integral. Its constant coefficient proves the last assertion.
\end{proof}

There is also a conditional-expectation interpretation. Draw the terminal
field with law $\rho_N\,d\lambda_N/Z_N$ and then follow the selected
block kernels. The conditional expectation of $\prod_{i\in S}F_i$
given the scale-$m$ field is
$[z_S]q_m/[z_\varnothing]q_m$ almost everywhere where the denominator
is positive. This follows by testing against a bounded function of the
scale-$m$ field. Its value on the zero-density set is immaterial.
Two histories generally have different conditional laws; the proposition
asserts equality of their integrated moments.

\section{Joint local terms and large-field components}

At an intermediate scale a local representation has the form
\begin{equation}
 e^{\Phi(U)+\sum_{S\ne\varnothing}z_S h_S(U)}
 \sum_{\mathcal F\ {\rm eligible}}
       \prod_{Y\in\mathcal F}
          \left(K_Y(U)+\sum_{S\ne\varnothing}z_S k_{Y,S}(U)\right).
 \label{ch11:local-source}
\end{equation}
Each regular term carries its required regularity condition;
each large-field term retains its original eligibility and large-field conditions.
An index $S$ may contain several insertion indices. Such terms are
generated by integration even if the regular source at the ultraviolet cutoff has only
terms linear in the source variables, that is, of \emph{source degree} one. The dependence set includes the link supports of the
insertions, event indicators, paths, and Gaussian windows. It need not coincide
with the exclusion support of a large-field activity.

One must never divide a coefficient with insertions by $K_Y$: that activity may
vanish. For example, the large-field parts of two insertion factors on a selected
large-field support contribute
\[
 K_Y(1+z_1b_1)(1+z_2b_2)
 =K_Y+z_1K_Yb_1+z_2K_Yb_2+z_1z_2K_Yb_1b_2.
\]
The joint coefficient contains the original activity once.

More generally expand the terminal product into its small-field
parts and large-field parts. Merge the link support of each large-field part with
every selected large-field support it meets, and merge intersecting
link supports of large-field parts. The exact cover ensures that the link support of a large-field part
containing a failure of the regularity condition meets at least one selected support.
Merge each resulting component, retaining the constituent indices,
original eligibility criteria, and union of their prescribed large-field sets.
Its value is the product of those constituents, with each selected
activity included once. Expanding the constituent indices recovers
the original term uniquely; hence this is an exact coefficient identity.
This algebraic construction gives all joint terminal large-field terms,
including at zeros of the original activities.

For a bound uniform in the number of inserted points, assume at most one
insertion point per site and a fixed finite list of link-support radii and
degrees. Then a site belongs to only a bounded number of link supports of
insertions. To apply the component tree estimate of Chapter~\ref{ch08:chapter}
to the union of the link supports, use a larger support exponent in the
bounds on the individual factors.
At the ultraviolet cutoff, with the Wilson action, every new component contains an original
large-field support, whose decay factor is
$\exp(-c\beta_Nt_N^2 n)$. Keep half this exponent aside before summing
attachments. A polynomial in $B_N,\widehat\beta_N$ times the remaining
decay factor tends to zero faster than every inverse power.
This proves small joint terminal bounds when the weighted large-field
parts of the insertions themselves have a summable site bound. That latter bound
is a quantitative hypothesis on the local section; it is not implied by the
finite polynomial identity.

\section{Support bounds for several components with insertions}

Fix $n$ insertion indices in each reflected replica cell. Suppose every
elementary principal term has at least one insertion point, contains its
insertion points in its dependence set, and that set is within
distance $\rho$ of each such point. Suppose insertion points in different
cells are separated by more than $2\rho$.
A component formed by intersecting these dependence sets then lies in one
cell. Indeed insertion points on adjacent terms have distance at most $2\rho$,
so they belong to the same cell. Connect the distinct insertion indices used by
the component. Erasing cycles gives a path of at most $n-1$ edges.
All its insertion points are therefore within $2n\rho$ of one another,
and every point of its dependence set is within $(2n+1)\rho$ of each
insertion point. This bound is independent of the number of elementary
terms in the component.

A second estimate controls increasing the source weights. Let each
component carry a nonempty index set $S$, dependence set $X$, and let
$t(S)$ be the number of replica cells represented in $S$.
Suppose at most $n$ indices occur per cell. Increase the common source
weight from $\omega_-$ to $\omega=H^U\omega_-$.
For components within one cell the factor is at most $H^{Un}$.
For components meeting several cells, retain an additional margin in the support exponent
$\sigma>0$. If
\begin{equation}
 \sigma c_B(X)\ge Un\,t(S)\log H,
 \label{ch11:spare-support}
\end{equation}
then the increased source weight and decreased support weight give
\[
 H^{U|S|}e^{-\sigma c_B(X)}\le1.
\]
Consequently an incidence bound $p_1$ at insertion points for single-cell terms
at exponent $e$, and $p_2$ for other terms at exponent $e+\sigma$,
give the bound $H^{Un}p_1+p_2$ at exponent $e$.

The geometric lower bound \eqref{ch11:spare-support} and the stronger
incoming norm are both required. An estimate proved only at exponent
$e$ cannot be reused at $e+\sigma$. Large-field components require
their own joint attachment estimate because one such component may
carry insertion points from many cells. The half of the decay exponent kept aside above, together
with the stronger support norm, absorbs this additional weight.
The support exponent is therefore chosen before the expansion with insertions
and before the history thresholds are chosen.

\section{One step with insertions}
\label{ch11:marked-step}

We give quantitative bounds before applying the Gaussian comparison.
This also specifies the margin in the support exponent needed when the number of
inserted points grows with the volume. Fix an output exponent $e\ge0$,
put $v=e+5$, and let $B_f\le2B_c$ be the adjacent support radii.
For the prescribed enlargement $X^\uparrow$ of a fine support, the
covering argument of Section~\ref{ch06:primitive-bounds} gives
\[
 Lc_{B_c}(X^\uparrow)\le15^4\bigl(L+10c_{B_f}(X)\bigr).
\]
Consequently, if $10\cdot15^4(e+7)\le L\sigma_f$,
\begin{equation}
 e^{(v+2)c_{B_c}(X^\uparrow)}
 \le T_e e^{\sigma_f c_{B_f}(X)},\qquad
 T_e=\exp(15^4(e+7)).
 \label{ch11:cover-transfer}
\end{equation}
The numerical constant describes this choice of support enlargement;
the essential feature is its fixed additive cost and $L^{-1}$ fine
cover cost.

In one sector, which includes one forest (one rooted spanning tree in every block), and for one value of the auxiliary index, distinguish three kinds of factors:
local corrections without insertions, incoming large-field factors, and insertion factors
with nonempty index sets. Their dependence sets determine
the intersection graph. The exact-cover condition requires the local corrections without insertions in each
component to cover every site prescribed by its incoming large-field factors. In particular a component
with no insertion-free correction has no incoming large-field factor.
This is a condition on the exact cover decomposition, retained when
factors with insertions are attached, rather than a consequence of spatial
intersection alone.

Let $\delta$ be the containing-site bound for insertion-free correction
amplitudes. Suppose the incoming large-field amplitudes, already
weighted by $e^{\sigma_f c_{B_f}}$, have site bound $\eta$.
Let $a$ bound the insertion incidence sum at a fixed insertion point after
\eqref{ch11:cover-transfer}, at exponent $v+2$.
If at most $g$ insertion points project to a coarse site, set
\[
 {\cal W}(B)=(33B)^4\sum_{j\ge0}(j+1)^4e^{-j},\qquad
 \Sigma=\delta+{\cal W}(B_c)L^4T_e\eta+{\cal W}(B_c)g a.
\]
The insertion and background incidence bound at a specified insertion point is
$T_e\eta+a$; no $L^4$ occurs there. The site bound is $\Sigma$ at
exponent $v+1$. To see the latter conversion, root a support at one of
its fine sites or insertion points. Its covering number is at least
$\lfloor d(o,x)/(8B_c)\rfloor$. The number of possible roots in shell
$j$ is at most $(33B_c)^4(j+1)^4$, while losing one support exponent
provides $e^{-j}$. Background roots have $L^4$ possible positions within the block;
roots at insertion points have at most $g$ choices.

Choose inflation factors $m,l\ge1$. The common local integral bound of
Chapter~\ref{ch06:large-fields} may be applied with the insertion-free
amplitudes inflated by $m$ for factors with a nonempty prescribed
large-field set, and by $l$ for other corrections. Smallness remains available
by increasing the lower bound on the reference coupling at fixed $m,l$.
A nonprincipal insertion factor is either pointwise at most its bounded
amplitude divided by $l$, or supported on an event of conditional
probability at most $l^{-1}$ and bounded by that amplitude.
For a compatible connected collection $C$, therefore,
\begin{align}
 t_C\left|\int\prod_{j\in C}F_j\,d\nu\right|
 &\le\prod_{j\in C}b_j,\notag\\
 t_C&=\begin{cases}
 m,&\text{some factor prescribes large-field sites},\\
 l,&\text{no sites are prescribed and some factor is nonprincipal},\\
 1,&\text{purely principal}.
 \end{cases}
 \label{ch11:integrated-component}
\end{align}
Here $\nu$ is the single Gaussian/Haar law of the sector, formed with its forest; by Hypothesis~(H) below the bounds $b_j$ and the factors $t_C$ do not depend on the forest.
If an insertion-free correction occurs, its inflated common product bound
absorbs $t_C$, while the other factors are bounded pointwise.
If none occurs, the condition on prescribed large-field sets removes incoming large-field
factors; use either the pointwise or event bound for one nonprincipal
insertion. No independence is used in either case.

Assume $(2B_c+1)^4\Sigma\le1$. The rooted-tree bound gives connected
component site sum at most $\Sigma$ at exponent $v$, and, at insertion points,
incidence at most $T_e\eta+a$. Indeed let
$c_i=c_{B_c}(X_i)$, $a_i=b_i e^{vc_i}$ and
$\widehat b_i=b_i e^{(v+1)c_i}$. The sum of $\widehat b_j$ over
regions meeting $X_i$ is at most $|X_i|\Sigma\le c_i$.
Iterating $R_i^{(0)}=a_i$ and
$R_i^{(q+1)}=a_i\exp(\sum_{j:X_j\cap X_i\ne\varnothing}R_j^{(q)})$
gives $R_i^{(q)}\le\widehat b_i$. Every finite connected collection
has a spanning tree and is included in this upper bound.
This is the rooted-tree majorization underlying the polymer criterion
of Chapter~\ref{ch08:chapter}; it is applied here with the full
dependence sets of terms with insertions, rather than proved anew for each subset of insertions.

The split of an incoming principal factor $Wz_S$ by a local cutoff is
exact:
\[
 (1+\gamma Wz_S)(1+(1-\gamma)Wz_S)=1+Wz_S .
\]
The second term is nonprincipal and the joint tail bound for cutoff remainders gives
its supporting event probability at most $x^{-1}$.
Thus $l\le x$ absorbs that term. If principal and arbitrary remainder
incidence bounds before enlargement are $h$ and $z$, respectively,
one may take
\begin{equation}
 a=2T_eh+lT_ez.
 \label{ch11:marked-incidence}
\end{equation}
Both pieces retain the same incoming event indicator and dependence set.
Purely principal connected components have the one-cell support
bound proved earlier. Their integral is the joint moment
in Proposition~\ref{ch11:joint-gaussian}, since it contains neither
an insertion-free correction nor an incoming large-field factor.

Finally normalize the insertion-free local gas by the same exact
connected expansion as in Chapter~\ref{ch09:renormalization}, that is, with its activities averaged over the forests as in Proposition~\ref{ch09:forest-average-prop}; the normalization constant does not depend on the forest (part~(i) of that proposition). Assumption~\ref{ch09:forest-assumption} concerns only the construction without insertions, so the combination of the fixed-forest computation above with this average requires the following hypothesis, which is not proved in this monograph. \emph{Hypothesis~(H) (forests with insertions).} For every forest $F$, the exact identity for the block integral of the incoming representation with insertions \eqref{ch11:local-source} has the form \eqref{ch09:fixed-forest-gas}, with the same constant $\mathscr N_F$ and function $\mathscr S$, with labels, supports $X_M$ (containing the dependence sets of the insertions carried by $M$) and prescribed large-field sets $I_M$ independent of $F$, and with activities $a_M^F(z,U)$ in $\mathcal A_{\mathcal I}$, Borel and gauge-invariant in $U$ and local on $X_M$, satisfying $\sum_M\sup_U\|a_M^F(\cdot,U)\|_\omega<\infty$; the labels comprise those of \eqref{ch09:fixed-forest-gas}, with the same supports and prescribed large-field sets and with coefficient of $z_\varnothing$ equal to the activity there, and further labels, which carry an insertion factor and whose coefficient of $z_\varnothing$ is zero; every coefficient $[z_S]a_M^F$ satisfies part~(b) of Assumption~\ref{ch09:forest-assumption}; the bounds of this section, in particular \eqref{ch11:integrated-component}, \eqref{ch11:marked-incidence} and the bounds $\Sigma$, $T_e\eta+a$, $k$, $\psi$, and every other estimate of this chapter which is made for one forest, in particular those of Section~\ref{ch11:actual-reader}, hold for every $F$ with right sides independent of $F$; and, for every $\varsigma\in\mathfrak G$, there is a bijection $M\mapsto\varsigma\cdot M$ from the labels for the insertions $F_\bullet$ onto those for the insertions $F_\bullet\circ\varsigma^*$, transporting supports and prescribed large-field sets as in part~(d), with $a^{\varsigma F}_{\varsigma\cdot M}[F_\bullet\circ\varsigma^*](z,U)=a^F_M[F_\bullet](z,\varsigma^*U)$, where $[F_\bullet]$ indicates the insertions from which the incoming representation with insertions is built. Under~(H), the proof of Proposition~\ref{ch09:forest-average-prop}(i)--(ii) applies coefficient by coefficient, because $\mathcal A_{\mathcal I}$ is commutative and the coefficients of activities with disjoint supports depend on the trees of disjoint sets of blocks: the identity holds with the averaged activities $\mathbb E_Fa_M^F$, which obey the same bounds, and its coefficient of $z_\varnothing$ is \eqref{ch09:averaged-gas}. Thus the density with insertions and the density without insertions are averaged by the same average over the forests. The estimates of this section are made for one forest and are used, through their uniformity in the forest, for the averaged activities.
If its local normalization corrections have total site bound $k$
at exponent $v$ after the $l$ inflation, choose an output bound
$\psi$ with $e^\Sigma-1+k\le\psi$.
The regular nonprincipal amplitudes are then bounded by $l$, the selected
components by $m$, and the incidence at insertion points by $T_e\eta+a$.
The corrections with empty index set contribute zero to that incidence.
In particular output incidence at most one follows from
$T_e\eta+a\le1$; it is not an extra independence or finite-arity
assumption. The estimates hold for every finite index set with
the stated bound $g$ on the number of insertion points per coarse site, uniformly in the number of replica cells.

\subsection{The one-insertion graph estimates}
\label{ch11:one-graph}

Fix a degree $d$ with coarse derivative order at least
$\max(2,d+1)$, fine spatial exponent at least $d$, fine smoothing order
at least $d+2$, and coarse smoothing order at least $d+3$ for the
simultaneous calculation for normalized local insertions and the regular and large-field remainders.
Assume the period is at least $100(1+B_c)$.
The fixed normalized local insertion is tested with one additional
derivative beyond each requested derivative used in a coefficient evaluation; this requirement
does not impose an additional derivative on an arbitrary incoming
activity.

Write $\kappa_i=(\epsilon_i^{\rm for})^{1/2}$ for the square root
of the large-field bound at scale index $i$.
The remainder norm used here is the sum of the anchored,
support-weighted regular majorants and $\kappa_i^{-1}$ times the
anchored, support-weighted large-field majorants, with their original
indices of support and prescribed large-field set. Thus the fine and coarse remainder norms
use $\kappa_f$ and $\kappa_c$, respectively, not a common value.

For this degree write $w_i$ and $f_i$ for the scaling
dimension and curvature-factor count of a homogeneous basis vector.
The leading rescaling is
$(S_Lx)_i=L^{d-w_i}x_i$, so $\|S_L^{-1}\|\le1$.
Taylor localization, the extraction of the local polynomial part by Taylor expansion of the fine local section on the free interpolated
coarse jet, gives the one-step localization map $T$, a triangular change of basis with diagonal
$(t_c/t_f)^{f_i}$. Its entries and inverse are bounded at fixed
degree and stage: the differentiated interpolation kernel has bound
$C_{d,r}L^{-1-r}e^{-c|y-y'|}$ for $r\le d$, summable against the
polynomial growth of a radial jet. Thus the leading map is $TS_L$.

Let $G$ be an incoming linear graph of regular vanishing-jet and
large-field remainders, with norm at most $s_f$ per coefficient. As in the graph transform for invariant manifolds, the step with insertions maps coefficient graphs at the fine scale to graphs at the coarse scale and halves the distance between two of them.
The step with insertions gives a linear coefficient map $M_G$ and a
raw remainder map $R_G$. For a common error bound $\epsilon$ above the bound $C_d/L+o_x(1)$, which is valid before $L$ and the lower bound on the reference coupling are chosen ($\epsilon$ may be taken to be $L^{-2/5}$ after those choices), it satisfies
\begin{align}
 |M_Gx-TS_Lx|&\le\epsilon(|x^\circ|+\|Gx\|_{\rm rem}),\notag\\
 \|R_Gx\|_{\rm rem}&\le C_d|S_Lx^\circ|
                                  +\epsilon\|Gx\|_{\rm rem}.
 \label{ch11:graph-estimates}
\end{align}
The circle deletes the identity-field coordinate. Finite symmetry averaging
only enlarges $C_d,\epsilon$ by fixed degree-dependent constants.

Here is the smallness argument for the remainder, including large
supports. Root the connected expansion at its unique factor with insertions;
the remaining factors have the common bound already established.
For a regular remainder on support $X$, put $r=c_{B_f}(X)$.
If its coarse cover exceeds $B_c$, one retained unit of exponential support decay gives
$E_L(\widehat\beta_f)e^{-B_c/2}$, where $E_L(x)\le C_Lx^{a_L}$ includes the Gaussian-scale costs and fixed derivative polynomials. This still tends to zero because $B_c\asymp(\log x)^2$.
Otherwise distinguish $r^3<L$ and $r^3\ge L$.
In the first case the support fits in the fixed centered nonwrapping
lift. Vanishing radial holonomy derivatives imply cancellation of screened
polynomial link moments: if direction $v_j$ is polynomial of degree
at most $a_j-1$ on the tested links and $\sum_j a_j\le d$, then
$D^k(R\circ\exp)(0)[v_1,\ldots,v_k]=0$.
To see this, write the local analytic coordinates as
$\log(e^{q(x)}H(J)_e e^{-q(x+e_\alpha)})$, with a site gauge
polynomial $q$ and radial coefficients $J$.
Their differential maps parameters of scaling dimension $a$ onto polynomial link
fields of degree at most $a-1$: decompose a homogeneous polynomial
one-form into its radial part and the gradient of
$(x\cdot A(x))/(a)$, then correct the lower-degree link-average error.
Higher derivatives preserve the summed scaling dimension.
Gauge invariance removes $q$; polarization of the vanishing holonomy derivatives and induction in the chain rule remove all partitions except
the desired derivative. Locality permits screening to the tested
links, always in the same centered lift.
The discrete Taylor decomposition of the actual interpolating field
therefore cancels its polynomial part through degree $d$; its
remainder contributes $L^{-(d+1)}$.
In the second case use
$1\le L^{-(d+1)}(1+r)^{3(d+1)}$.
The additional incoming exponential support decay absorbs this polynomial and
all derivative costs. Multiplication by the insertion rescaling $L^d$
therefore leaves $C_d/L$ in both cases.
For a large-field input of physical support-weighted mass $m$,
the incoming normalized mass is $m/\kappa_f$. The large-field factor carrying this input lies in no purely
principal component. The nonprincipal and selected-component estimates,
followed by local subtraction, give
\[
 C_{\rm orb}|x_{\rm out}|\le\eta m/\kappa_f,\qquad
 \|R_{\rm out}\|_{{\rm rem},\kappa_c}\le\eta m/\kappa_f,
 \qquad \eta=L^{-2/5}.
\]
The underlying real/imaginary input decomposition contributes $2m$, not
$2m/\kappa_f$. The factor $\kappa_f^{-1}$ enters only the input norm;
the outgoing selected terms are weighted by $\kappa_c^{-1}$ in the
local subtraction. The common large-field energy bound absorbs that
outgoing normalization before the displayed estimate is formed.
Both adjacent bounds are retained in this comparison.
A normalized local insertion contributes its extracted polynomial and a bounded
Taylor remainder; its coefficient error has the same small factor.
These are the three terms in \eqref{ch11:graph-estimates}.
A wrapping small-field remainder is never declared large-field.

Choose $L$ and then the lower bound on the reference coupling so the coefficient
perturbation is smaller than half the inverse bound of $T$.
Absorption gives $|x|\le s_f|M_Gx|$ after fixing a suitable stage
column constant $s_f$. Thus $M_G$ is invertible.
The outgoing graph $G'=R_GM_G^{-1}$ lies in the prescribed ball
$\|G'\|\le s_c$ if $s_c$ is chosen above the finite constant from
\eqref{ch11:graph-estimates}. This ball invariance is part of the
choice, not a consequence of invertibility alone.
The identity field has zero remainders and is sent exactly to $L^d$ times itself.

All arrays and maps taking local polynomial coefficients are fixed independently of $G$.
For graphs $G,H$ with $a=\|G-H\|/s_f$, the distinguished factor in
their difference is $(G-H)x$; the normalized local insertion cancels.
The same rooted sums therefore prove
\begin{equation}
 |(M_G-M_H)x|+\|(R_G-R_H)x\|_{\rm rem}
                         \le\delta a|x|,
 \qquad \delta=C_d\epsilon s_f .
 \label{ch11:graph-difference}
\end{equation}
The harmless $C_d$ covers the finite symmetrization.
Increase the block ratio and the lower bound on the reference coupling until $4\delta s_f\le1$,
together with the preceding inverse and invariance inequalities.
These are raw-map difference estimates; continuity alone
would not supply them. They are the input for the contraction estimate
of the graph transform on the coefficient graphs.

For fixed incoming data the maps are continuous in the inverse coupling
and the canonical path coefficients in weighted $\ell^1$ with a strict margin in the decay width, and in the incoming activity arrays in the same strong norm. All local integrands use the same
protected physical-coordinate change and fixed common index sets.
The common integrable bounds in Chapter~\ref{ch06:large-fields}
justify dominated convergence on each fixed finite truncation. The strict margins in the path width and in the support exponent make the omitted countable tails uniformly small. Taking the finite truncation first and then its size to infinity proves continuity in these weighted norms. Finite Taylor localization and symmetry
averaging preserve continuity. Inversion preserves it on the common
inverse-norm ball. No extra derivative of an incoming activity is
taken when comparing equal-regularity adjacent stages.

\subsection{Propagating the already formed terms of source degree at least two}
\label{ch11:higher-modes}

Suppose every pair of distinct ancestors with insertions has periodic
maximum-distance at least $B^2$. If a dependence set contains $r\ge2$
such ancestors, then
\begin{equation}
 Br\le64c_B(X).
 \label{ch11:separated-cover}
\end{equation}
A doubled spanning traversal of a minimal connected cover has length
at most $6B(c_B(X)-1)$. Assign each ancestor to a covering center.
The assigned centers are distinct and separated by at least $B^2-2B$.
Their cyclic traversal costs at least $r(B^2-2B)$. For $B\ge8$
the stated constant follows. This argument uses periodic distances.

Retain two cover units beyond those needed by the connected expansion.
If the source weights grow by $\widehat\beta_f$ per insertion, then eventually
\[
 \widehat\beta_f^{|I|}e^{-c_B(X)}
       \le e^{-c_B(X)/2},\qquad |I|\ge2,
\]
because $64\log\widehat\beta_f/B\le1/2$.
The second unit absorbs site conversion and attachments. Existing terms of source degree at least two thus gain at least $e^{-B/64}$ while retaining their final support
weights. This uses the full incoming majorant for terms of source degree at least two before
changing weights; the single-insertion bound alone would not give it.

Let $H_f$ be that stronger incoming incidence bound and $\zeta$ the
full weighted incidence sum for terms of source degree one. With a strict margin in the
site criterion for insertion-free factors of Section~\ref{ch11:marked-step}, root each
component at an existing factor of source degree at least two or at a factor of source degree one
in a newly merged component. The common integral and rooted-tree estimates give
\begin{equation}
 H_c\le C_e\bigl(e^{-B/64}H_f+\zeta\bigr).
 \label{ch11:higher-induction}
\end{equation}
The constant includes the fixed change of cover radius and the finite list
of degrees, but not the period or number of insertions.
A new component of source degree at least two has at least two individual insertion indices; using the bound for
one gives the deliberately weaker bound by $\zeta$.
Existing large-field coefficients of source degree at least two are treated as whole incoming factors
with their prescribed large-field sets retained. Newly selected components use the
margin in the decay exponent of Chapter~\ref{ch06:large-fields} to absorb the fixed
inflations in \eqref{ch11:integrated-component}.
Choose the stage threshold so $C_e e^{-B/64}\le1/2$ and
$C_e\zeta\le1/2$. Then $H_f\le1$ implies $H_c\le1$.
The joint terminal construction starts the induction.

If the input is specified at exponent $\sigma_f-4$, its change
of cover radius must still leave the two additional coarse units. For example,
\[
 10\cdot15^4(e+9)\le L(\sigma_f-4)
\]
suffices for the enlargement used here.
The four removed fine units and the two retained coarse units are
different costs. A declaration that exponents lie in $[98,100]$
would not prove this inequality.

\section{A joint Gaussian comparison}

The comparison of several local polynomials must preserve their cross
covariances. Separate Gaussian approximations for individual factors
are insufficient. The following finite-degree statement is the useful
form, permitting a different cutoff for every subset of insertions.

\begin{proposition}[Joint polynomial comparison with cutoffs]
\label{ch11:joint-gaussian}
Fix $n,D$ and $A,B\ge1$. For each nonempty $I\subseteq\{1,\ldots,n\}$,
let $p_I$ be a gauge-invariant polynomial of degree at most $D$ and
coefficient $\ell^1$ norm at most $A$.
For every member $t$ of a finite probability mixture, let one probability
space carry a centered joint Gaussian array $\eta_t$, common to all
subsets, with covariance $C_t$. Suppose
\[
 |a_j|\le B,\quad |(C_t)_{ij}|,|(C_*)_{ij}|\le B,\quad
 |(C_t-C_*)_{ij}|\le\varepsilon,\qquad 0\le\varepsilon\le1.
\]
For each $I$ let $0\le\gamma_{t,I}\le1$ satisfy
$\mathbb E_t(1-\gamma_{t,I})\le\varepsilon^2$.
On its positive set suppose the actual normalized local curvature
coordinates, after an allowed gauge rotation, obey
$|\xi_{t,I,j}-a_j-\eta_{t,j}|\le\varepsilon$ on all used coordinates.
The polynomial's invariance removes the rotation.

With the same probability weights $\lambda_t$ for all subsets and the empty product,
define
\[
 M_S=\sum_{\pi\in\mathcal P(S)}\sum_t\lambda_t
       \mathbb E_t\prod_{I\in\pi}\gamma_{t,I}p_I(\xi_{t,I}),\qquad
 M_S^*=\sum_{\pi\in\mathcal P(S)}
       \mathbb E_{C_*}\prod_{I\in\pi}p_I(a+\eta),
\]
and set $M_\varnothing=M_\varnothing^*=1$.
The symbol $\mathbb E_{C_*}$ may denote the Gaussian pairing
functional for a symmetric array, without assuming positivity.
There is a polynomial $P_{n,D}$ such that
\[
 \sum_{S\ne\varnothing}|M_S-M_S^*|
       +\|\log M-\log M^*\|_1
       \le\varepsilon P_{n,D}(A,B).
\]
Its constants are independent of the number of unused coordinates
and of the number of mixture members.
\end{proposition}
\begin{proof}
Fix a partition with $\ell$ blocks and put
$\gamma=\prod_{I\in\pi}\gamma_{t,I}$.
Then $1-\gamma\le\sum_I(1-\gamma_{t,I})$, so
$\mathbb E_t(1-\gamma)\le\ell\varepsilon^2$ without independence.
Telescope each product monomial to replace the physical coordinates
by $a+\eta_t$. Every difference contains one factor of size
$\varepsilon$ and at most $\ell D-1$ Gaussian factors.
Fixed marginal moments and H{\"o}lder's inequality bound the result by
$\varepsilon A^\ell C_{n,D}(1+B)^{nD}$.
Only coordinates in that monomial occur; no ambient-dimension count
is introduced.

Cauchy--Schwarz removes $\gamma$ at cost
$\sqrt\ell\,\varepsilon$ times the second moment of the polynomial.
Finally the Gaussian pairing formula~\cite{Isserlis1918} expresses each
reference moment as a finite sum of products of means and covariances.
Telescope these products to replace $C_t$ by $C_*$.
This contributes another $\varepsilon$ times a fixed polynomial in $A,B$.
Average with the common probability weights and sum over finitely many
partitions. The finite polynomial formula \eqref{ch11:cumulants} gives
the logarithm bound. The empty coefficient is exactly one because
the empty product has no cutoff.
\end{proof}

The proposition applies to the purely principal components of an
expansion with insertions, after the normalization constant has been extracted once.
It does not identify the whole interacting density with a Gaussian.
Deleting an additive remainder for an estimate may leave the precision,
affine mean and prescribed path data unchanged, but it changes the
interacting density. Exact reconstruction always uses the full
density. This distinction is indispensable when matching the Gaussian
comparison to the step.

\subsection{The exact cross-covariance in the block axial gauge}
\label{ch11:tree-covariance}

We identify the flat reference used in the preceding comparison.
Use coordinates with the unnormalized $\ell^2$ scalar products on the two adjacent tori.
Let $D,D'$ be their bond-to-plaquette curls and let
\[
 K=D^*M_fD,\quad K'=(D')^*M_cD',\quad
 A=K+bQ^*Q ,
\]
where $M_f,M_c$ are the modified positive curvature metrics and $Q$
is the rooted linearization of the averaging operation. Let $I_F$ insert the non-tree
coordinates and $A_F=I_F^*AI_F$. The tree gap gives $A_F>0$.
Set $X=I_FA_F^{-1}I_F^*$.
The rooted gauge projection $\mathcal P$ onto the slice of the same block axial gauge
satisfies $A\mathcal P=A$ and $D\mathcal P=D$, since its gauge
potential vanishes at each block root. Thus $XA=\mathcal P$
and $DXA=D$; no inverse of the gauge-degenerate full $K$ is used.

Choose matrices $Z,Z'$ solving $KZ=D^*$ and $K'Z'=(D')^*$.
They exist because $\ker K=\ker D$ and the symmetric matrix $K$
has range $(\ker D)^\perp$; similarly for $K'$.
The matrices $P_f=DZ$ and $P_c=D'Z'$ are independent of the choices,
symmetric and positive semidefinite.

For the linearization $W$ of the critical configuration, taken in its representative $W_F$ in the block axial gauge, the adjacent
Gaussian identities of Chapter~\ref{ch05:gaussian} read
\[
 KW=Q^*K',\quad AW=bQ^*,\quad DW=JD',
\]
where $J=H_{0,f}^{-1}T H_{0,c}$ is the curvature response.
Here $T$ is the induced map on curvatures and $H_{0,f},H_{0,c}$
are the unmodified diagonal operators; the modified metrics
$M_f,M_c$ still define $K,K'$ and agree with them on curls.
Thus $J$ is not the induced map $T$.
Transpose the first identity and multiply by $Z$:
$K'QZ=(D')^*J^*$. Multiplying by $(Z')^*$ gives
$D'QZ=P_cJ^*$. On the other hand $AZ=D^*+bQ^*QZ$.
Multiply by $DX$ and use $DXA=D$ and $DXbQ^*=DW$ to obtain
\begin{equation}
 (DI_F)A_F^{-1}(DI_F)^*
       =DXD^*=P_f-JP_cJ^*.
 \label{ch11:conditional-curvature}
\end{equation}
This proves every cross-entry simultaneously, including different
plaquettes and coordinates in the basis of Section~\ref{ch01:gauge}. It retains both modified metrics,
even when an adjacent pair straddles a change of stage.
The right side is positive because the left is a covariance.
Harmonic and gauge modes were handled through the range equations,
not discarded.

At zero background the Gram matrix of the components is precisely $A_F$:
the absolutely summed path coefficients reproduce their curvature
factorization, the linearization of the averaging operation is $Q$, and each signed
orientation has the same quadratic form. Their probability average
therefore adds no extra factor. The path-length truncation and the non-flat
background are handled by the inverse-window comparison, not by
changing this exact identity.

\subsection{Nonlinear local coordinate and cutoff estimates}
\label{ch11:actual-reader}

Fix a clean neighborhood and one choice of spanning trees of the blocks (a forest; the identity \eqref{ch11:conditional-curvature} holds for every forest and its right side does not depend on the forest, and the estimates below are used, by Hypothesis~(H) of Section~\ref{ch11:marked-step}, with bounds independent of the forest, so that they pass to the average over the forests), one value of the auxiliary index, one coarse gauge and one Gaussian
sample for all coordinates read by the polynomial family.
Section~\ref{ch06:linear-comparison} gives, against the flat model, mean error
$O(B^a t_f^2)$ and covariance error $O(B^a t_f)$, plus exponentially
small window tails. After multiplication by the finitely many curl
rows and summing the finite-range window entries, these estimates
remain $x^{-1/2+o(1)}$.
Combining with \eqref{ch11:conditional-curvature} gives the full
conditional curvature array, rather than only its diagonal.
All comparisons use the same local gauge transformation based at the root; the gauge-invariance of each
polynomial then permits its physical evaluation with a different local
gauge transformation without changing its value.

Here is an explicit scale choice that makes the nonlinear error small.
Put $R=x^{-3/8}$ and $g=\beta^{-1/2}$ on the coupling band.
The coarse logarithm on the required patch has size
$x^{-1/2+o(1)}=o(R)$.
Restrict the finitely many Gaussian read coordinates so the physical
fluctuation is at most a fixed fraction of $R$.
Since $R/g=x^{1/8+o(1)}$, the Gaussian tail bound makes failure of
this restriction smaller than every inverse power of $x$, even after
the polynomial-in-$B$ count of the local reads.
On this event, the product-logarithm and curvature Taylor remainders
are $O(B^a R^2)$. Dividing by $g$ gives
$O(B^a x^{-1/4})$.
The mean error against the flat model, in Gaussian units, is
$O(B^a t_f^2/g)=x^{-1/2+o(1)}$.
The omitted path tails are smaller than every power of $x$.

For a smooth cutoff also requiring the fine small-field cutoff, use the
conclusion of Proposition~\ref{ch06:mean-cancellation} for event indicators far from the large-field set.
The linear fine comparison component is $O(L^{-2}t_f)$; choose the fixed
block ratio before the lower bound on the reference coupling so it lies strictly within
the fine plateau. The comparison for the averaging operation has its independently
chosen strict plateau width. Restricting the local fluctuation variables to the smaller
fixed fraction of $t_f/g$ required by that proposition has failure
smaller than every power of $x$, because $t_f/g=M_f^{\rm cut}$ grows
faster than every power of $\log x$. Intersect these cutoff events
with the radius-$R$ event. Their failure bound is still summable over
the finitely many local reads.

Thus for every fixed family of polynomial reads with the prescribed
clearance and support, the hypotheses of
Proposition~\ref{ch11:joint-gaussian} hold with, for example,
$\varepsilon=x^{-1/32}$ after increasing the lower bound on the reference coupling.
The covariance error is smaller than $\varepsilon$ and cutoff loss
smaller than $\varepsilon^2$. One joint Gaussian sample and one linear prediction of the fine
curvature from the coarse field are used for all subsets. The normalization constant, which does not depend on the forest (Proposition~\ref{ch09:forest-average-prop}(i)), is the one
extracted from the integral without insertions before the insertions are chosen; all
remaining principal products are integrated against its normalized
Gaussian/Haar probability law. Incoming large-field and correction
factors belong to the other components already bounded in
\eqref{ch11:integrated-component}, and are not absorbed into a
subset-dependent probability measure.

\section{Restarting below a comparison scale}

A useful auxiliary history, the \emph{restarted history}, keeps the inverse coupling constant below a
comparison scale while retaining the paths, adjacent quadratic identities,
and block geometry. Proposition~\ref{ch11:mass} preserves all original
normalized coefficient integrals along this history. The construction
below provides its local representations and estimates on that same
restarted history.

Two changes of representation are involved. For a large reference coupling parameter
$s$, define
\[
 H_s=e^{(\log\log s)^2},\qquad h_s=e^{(\log\log s)^{3/2}},\qquad
 t=s^{-1/2}H_s,\quad \vartheta=s^{-1/2}h_s .
\]
The restarted large-field threshold is $\vartheta$, not $t$; its
support radius is the square of the original logarithmic-square radius.
Eventually $2\vartheta\le t$. For the piecewise quadratic--linear function
$d_t$, with linear coefficient $c_{\rm lin}<c_{\rm dir}/16$,
\[
 d_t(r)-d_\vartheta(r)\ge0,\qquad
 d_t(r)-d_\vartheta(r)\ge c_{\rm lin}\vartheta^2
                      \quad(r>\vartheta).
\]
For $\vartheta<r\le t$ this follows from
$c_{\rm dir}r^2/2-c_{\rm lin}\vartheta r$; for $r>t$ it is
$c_{\rm lin}(t-\vartheta)r$.
The positive difference controls newly selected large-field sites.
The next construction assigns their event indicators, eligibility criteria, joint
insertion terms and path tails consistently.

% !TEX root = ../monograph.tex
\subsection{Constructing the restarted representations}
\label{ch11:restart-construction}

We now construct the auxiliary history just described. The physical
insertions retain their original calibration; only the integration
below the comparison level is restarted. All support norms in this
subsection use connected covers and include full dependence regions.
Their exponents are chosen with the reserves required by the marked
expansion before choosing an inverse coupling threshold.

Write $x_i=\widehat\beta_i$ and, for $x$ large, put
\[
 H(x)=e^{(\log\log x)^2},\quad h(x)=e^{(\log\log x)^{3/2}},\quad
 b(x)=\lceil(\log x)^2\rceil,\quad
 \vartheta(x)=x^{-1/2}h(x).
\]
The new representation at level $i$ has radius $B_i=b(x_i)^2$
and threshold $\vartheta(x_i)$. Its clip has the original two
coefficients, evaluated at this new threshold. Its path coefficients
and block-kernel branches are the original ones. In particular,
``retaining the clip'' means retaining its coefficients and formula,
not asserting equality of two different threshold functions.

\paragraph{Choice of reference parameter for the cutoff.}
There is a unique sufficiently large $r(x)$ satisfying
\[
 r(x)^{-1/2}H(r(x))=\vartheta(x).
\]
Indeed the function on the left is eventually strictly decreasing;
its values at $x$ and $x^2$ bracket the right side. Thus
$x\le r(x)\le x^2$, and taking logarithms gives
$r(x)=x^{1+o(1)}$. With this choice the original threshold formula evaluated at $r(x)$
gives the new threshold at $x$. The inverse coupling in the block
integral still lies in the band determined by $x$. In particular
\[
 b(x)\le b(r(x))\le4b(x),\qquad
 25b(r(x))+30\le B_i\le b(r(x))^2
\]
eventually. On each fixed adjacent stage pair,
$x_{i+1}=x_i+\kappa_G\log L_i$ implies
$B_{i+1}\le2B_i$ and
$\vartheta(x_{i+1})/\vartheta(x_i)\to1$.
The local chart, clipped source, Gaussian precision and tail events
are those of Chapter~\ref{ch06:large-fields} with reference parameter $r(x_i)$
and threshold $\vartheta(x_i)$. Its elementary scale inequalities
continue to hold: the radii are logarithmic powers and the thresholds
are $x^{-1/2+o(1)}$. The actual inverse coupling band remains expressed in
$x_i$, not in $r(x_i)$. In particular one must not invoke a lower
bound $\beta\ge r(x_i)/2$. For every fixed $\eta>0$, changing the reference parameter and
the actual band instead give, eventually,
\[
 y^{1-\eta}\le\beta\le2y,\qquad y=r(x_i).
\]
For a prescribed finite local estimate, fix the number of powers
of $\beta^{-1}$ and $g=\beta^{-1/2}$ that occur before choosing
$\eta$. Replacing a comparison power of $y$ by the corresponding
actual inverse coupling power then loses at most $y^{C\eta}$, where $C$
is fixed by that estimate. Choose $\eta$ small enough to retain
its negative exponent. The dimensionless quadratic gap and path
bounds are unchanged; physical covariance factors retain their
actual $\beta^{-1}$ normalization. Likewise, the cutoff-to-noise
ratio is computed exactly as
$\vartheta(x_i)/g\asymp h(x_i)$, and the selected-component energy
is $\beta\vartheta(x_{i+1})^2\ge c h(x_{i+1})^2$.
These comparisons, not a fictitious inverse coupling band in $y$, are
used below.

\begin{proposition}[Re-expression of the endpoint]
\label{ch11:restart-endpoint}
At a sufficiently late comparison level $k$, the original density has an
exact presentation with radius $b(x_k)^2$, threshold
$\vartheta(x_k)$, the original path coefficients, clip coefficients,
inverse coupling and scalar, regular bound no larger than its original
bound, and large-field site bound
\[
 G(x_k)=\exp\bigl(-a_*h(x_k)^2\bigr).
\]
Here $a_*>0$ is chosen sufficiently small at each fixed stage before
its threshold for the starting scale. The construction retains the old large-field
labels and their eligibility criteria inside the new activities.
\end{proposition}
\begin{proof}
Abbreviate $s=x_k$, $t=s^{-1/2}H(s)$, $\vartheta=\vartheta(s)$,
$M=b(s)$ and $B=M^2$. Write $A_t,A_\vartheta$ for the path actions
with their respective event indicators and clips, and $E_t,E_\vartheta$ for
the original regular coefficients with the corresponding event indicators.
The old regular bound is at most $s^{-11/10}$; the old
large-field bound is at most $w_s=e^{-H(s)^{1/5}}$.
These are the selected trajectory's endpoint bounds.

Eventually $2\vartheta\le t$. The clip comparison proved above gives
$d_t(r)-d_\vartheta(r)\ge c_{\rm lin}\vartheta^2$ whenever
$r>\vartheta$, and a nonnegative difference otherwise.
Distribute each plaquette's difference among its incident sites.
Allowing five sites is a harmless upper bound. For the set
$F_\vartheta$ of new large-field sites this gives local energy contributions $d_z$ with
\[
 \sum_zd_z=\sum_p(d_t(r_p)-d_\vartheta(r_p)),\qquad
 d_z\ge c_{\rm lin}\vartheta^2/5\quad(z\in F_\vartheta).
\]
The omitted quadratic path rows give another nonnegative energy contribution.
Truncate each row at true path length $B$ before squaring, and keep
its half-square only where the old retention condition holds and the new one
fails. The relaxed small-field condition reads a radius-$B$ neighborhood: outside it
$2\le\vartheta e^{B/c}$ makes the condition automatic. The truncated
half-square therefore reads radius at most $2B$ and vanishes if its
tested neighborhood is new-small-field. Distribute its value equally
among the new large-field sites in that neighborhood, making no division
when the set is empty. The resulting local energy contributions $b'_z\ge0$ depend only on a neighborhood of radius
at most $4B$, and their sum is exactly the sum of truncated
half-squares. Perform this operation separately for each signed
orientation, retaining its averaging weight.

Choose $a_*\le c_{\rm lin}/20$, put $u=a_*h(s)^2$, and for large-field sites
$z$ put $X_z=\exp[-\beta(d_z+b'_z)]$. The actual band gives
$\beta\vartheta^2\ge h(s)^2/2$, hence $|X_z|\le e^{-2u}$.
The remaining exponent difference is
\[
 D=-(A_t-A_\vartheta)+\beta\sum_z(d_z+b'_z)
                               +(E_\vartheta-E_t).
\]
It consists of the omitted quadratic row pairs, the cubic, quartic
and loop event indicator differences, and the regular event indicator differences.
Each term vanishes on a new-small-field full dependence region.
At a fixed larger support exponent $a+2$ its site bound is
\begin{equation}
 \epsilon_s\le C B^C\bigl(
 \beta C_{\rm path}^2e^{-\sigma B/4}
 +\Gamma_s t^2+\beta C_{\rm path}(t^3+t^4)+s^{-11/10}\bigr)
                  =s^{-1/2+o(1)}.
 \label{ch11:restart-endpoint-error}
\end{equation}
Here $\sigma>0$ is a retained true-length path rate and
$\Gamma_s\le2H(s)$ is the loop coefficient bound.
For the first term expand the difference of squares into ordered
path pairs, one of length exceeding $B$, and retain part of their
exponential length weight. For the atoms restricted by event indicators use their old
bounds $t^j$ and the cover bound $1+2j+\operatorname{length}/B$.
For an old support $Y$, its fixed halo obeys
$c_B(Y^{[8]})\le C+C(M/B)c_M(Y)$. Eventually its old exponent
absorbs this new exponent and the extra unit used to turn an anchored
bound into a site bound. These are the support conversions of
Section~\ref{ch06:primitive-bounds}; their losses are polynomial in
$B$. The same argument bounds enlarged old large-field labels by
$CB^Cw_s$. No path multiplicity or winding path is removed.

Coalesce the absolutely summable terms by full support into $D_Y$.
For large $s$, put $W_Y=e^{D_Y}-1$, so $|W_Y|\le2|D_Y|$ and
$W_Y$ vanishes when its support is new-small-field. The identity
\begin{equation}
 e^{-A_t-E_t}=e^{-A_\vartheta-E_\vartheta}
                \prod_Y(1+W_Y)\prod_{z\in F_\vartheta}X_z
 \label{ch11:restart-endpoint-identity}
\end{equation}
is exact. Multiply by the old large-field gas and expand the optional
$W_Y$ factors. Keep the old compatibility and eligibility criteria.
Group using intersections of full dependence supports, including the fixed
halo for old large-field conditions and the radius-$4B$ supports of $X_z$.
Every nonzero component contains a selected $X_z$: an old activity
contains an old large-field site, hence a new large-field site, and every nonzero $W_Y$
reads a new large-field indicator. The factor at that site is mandatory and intersects
that component. Assign the component its exact new set of large-field sites and also
retain the union of its old sets of large-field sites and all their conditions.
Conversely, those constituents recover the original compatible
collection and optional subset. Thus the regrouping is a bijection
of terms, not a replacement of old eligibility by a new condition.
Disjoint new exclusion supports are the unions of these full read
supports. Their eligibility criteria require exact coverage of the new large-field sites there.

Reserve $e^{-u}$ per selected site factor. The remaining site sum
at exponent $a+1$ is at most
\[
 \rho_s=2\epsilon_s+CB^Cw_s+
       (8B+1)^4e^{(a+1)9^4}e^{-u}=s^{-1/2+o(1)}.
\]
A radius-$4B$ box has a connected $B$-cover of at most $9^4$
boxes; at most $(8B+1)^4$ such boxes contain a given site.
The connected-set bound of Chapter~\ref{ch08:chapter}, losing one
support exponent, applies since $(2B+1)^4\rho_s\le1$.
Every component retains at least one reserved factor, so its total
site bound is at most $e^{-u}(e^{\rho_s}-1)\le G(s)$.
There is no extra sum over prescribed sets of large-field sites: selected site factors already
specify them, and exact-coverage conditions only remove terms.
Absolute convergence justifies countable labels and coalescing.

The regular coefficients themselves are unchanged. Their domains
and allowed directions shrink with the threshold, and increasing
the cover radius decreases their support weights. Their old bounds
therefore persist. The exact identity and the gas bijection prove
the assertion, with unchanged scalar. All operations preserve gauge
invariance, reality and signed symmetry, by reindexing the original
orientation average.
\end{proof}

\paragraph{The genuine density and its neutral reference.}
At an intermediate level write the genuine presentation as
$R=e^{-E}R^0$, where $R^0$ is obtained by deleting only the regular
additive family $E=\sum_\alpha E_\alpha$. Its path, clip, large-field
activities and all their conditions remain unchanged. If the anchored
weighted value bound of that family is $q$, then
$\sum_\alpha|E_\alpha|\le q|\Lambda|$; this proves integrability
at finite volume, but is not used as a volume-uniform estimate.
Instead insert local factors
\[
 h_\alpha=e^{-E_\alpha}-1,\qquad
 \prod_\alpha(1+h_\alpha)=e^{-E}.
\]
Since $|h_\alpha|\le e^{|E_\alpha|}|E_\alpha|$, their local site
bound is the old anchored bound times a fixed polynomial in the
radius. They are inserted into the same expansion as the neutral
Gaussian integral. Thus the object integrated is always $R$, also
for every coefficient of a marked density. Deleting $E$ for the
neutral estimates never deletes it from the density identity.

\begin{proposition}[One restarted history on the comparison window]
\label{ch11:restart-family}
Fix the finite degree and derivative requests and the comparison
window with top level $k$, reference inverse coupling $s=x_k$, and at most
$(\log\log s)^4$ steps. Take its thresholds for the starting scales large enough
that $3s/4\le x_i\le s$ and
$|\beta_{N,p,k}-s|\le s/8$. On the original adjacent path data there
are presentations $P_i$, $i\le k$ in that window, with the bounds
$(B_i,\vartheta(x_i))$, common actual inverse coupling
$\beta=\beta_{N,p,k}$, large-field bound $G(x_i)$ and additive
value cost at most
\[
 q_{k-i}=s^{-1/200}A_s^{4(k-i)},\qquad
                    A_s=e^{(\log\log s)^3}.
\]
Their exact densities satisfy
\begin{equation}
 D_k=\rho_{N,p,k},\qquad
 D_i=\mathcal B_{i,\beta}D_{i+1}=\operatorname{density}(P_i),
 \qquad \int D_i=Z_{N,p}.
 \label{ch11:restart-density}
\end{equation}
Here the block operator has the original geometric branches, evaluated
at the common inverse coupling. The same local factors, normalization and
all deletion-of-support identities supply both weighted marked-step
uses. The constants are uniform in cutoff, period and the compatible
local fields, for the fixed requests.
\end{proposition}
\begin{proof}
Start with Proposition~\ref{ch11:restart-endpoint}.
The reference-inverse-coupling band gives $x_i/2\le\beta\le2x_i$ on every step.
Only the actual inverse coupling is frozen; the reference inverse couplings, original path pairs,
metric floors, gauge-covariant block average derivatives and forest charts keep their
original adjacent values. The recording-scale inequalities above
verify the radii and threshold comparisons in the local integral.
The fine path and clip in that integral are exactly the ones in the
incoming presentation. A change of stage uses its own adjacent pair,
not the previous stage's quadratic form.

Suppose the genuine incoming value cost is $q_n$, $n=k-i-1$.
Its neutral version has zero such cost and the same large-field
bound and loop coefficient bound. Thus it meets the small incoming
regular bound in Chapter~\ref{ch06:large-fields}; we do not claim
that the larger $q_n$ meets that derivative bound. The genuine
additive factors are handled by the local multipliers just defined.

We spell out the quantitative margins on this actual expansion.
The primitive exponent decomposition in
Section~\ref{ch06:primitive-bounds}, with reference parameter
$y=r(x_i)$, gives an unselected local exponent of site size
$t_y$ with $2t_y\le y^{-1/48}$ after enlargement of its threshold.
Indeed the corresponding exterior row, determinant and Haar
bounds are $y^{-1/4+o(1)}$, including radius polynomials and the
actual inverse coupling factors just discussed; the linear residual is
bounded by $y^{-4}y^2$ before these subpower losses. The two
determinant-word series and inverse correction are exponentially
small in $b(y)$, including the $y^4$ factor for two Gaussian
coordinates and any fixed inverse powers of the actual inverse coupling. The inherited path
terms have costs $O(\Gamma\vartheta^2)$ and
$O(\beta(\vartheta^3+\vartheta^4))$, with true-length reserves.
The inherited bounds before changing the reference parameter can be weakened to
$y^{-1/12}$, with fixed finite powers of the actual inverse coupling;
choose the preceding $\eta$ so their loss is less than $1/48$.
Thus all these are negative powers with a strict margin beyond
$1/24$. One may weaken their sum successively to $y^{-1/24}$
and then $y^{-1/48}$
to include exponentiation, support conversion and fixed constants.
The neutral additive part is zero. The local exponents are bounded
extensions on the same chart event indicators; outside them the failed cutoff conditions
are retained as tail-event factors, not evaluated as small exponents.

Selected components have, after the extracted Gaussian and local
normalization factors have been bounded, bound
$\exp[-K h(x_{i+1})^2 r]$ for $r$ selected cutoff conditions, by the actual
component energy in Chapter~\ref{ch06:large-fields}. Choose
$a_*>0$ also so small that $4a_*\le K/2$.
For both marked-step uses at once take
\[
 m=\max\{G(x_i)^{-2},G(x_{i+1})^{-1}\},\qquad
                 1\le l\le\max(1,H(s)^{q-3}),
\]
where $q$ is the fixed requested finite power.
The ratio $h(x_i)^2/h(x_{i+1})^2$ tends to one uniformly in the
window: its logarithm is $O(\sqrt{\log\log s}/\log s)$.
Thus $\log m\le4a_*h(x_{i+1})^2$ eventually, retaining at least
half the component energy. Apply the factor $m$ only to components with an
actual selected prescribed set of large-field sites, and to a selected component within each
such component. The original prescribed set of large-field sites coverage ensures that an
incoming large-field factor cannot occur without such a component.
Factors $l$, finite decorations and local counting costs have
logarithms $O((\log\log s)^2+\log B_i)$ per selected cutoff condition and fit
inside the remaining margin.

The simultaneous tail estimate of Chapter~\ref{ch06:large-fields}
uses these very chart coordinates and event indicators under the same Gaussian
law, including any bounded genuine multipliers. It therefore bounds
all compatible factors from failed cutoff conditions jointly. Its cutoff exponent
and the retained selected-component exponent dominate every multiple
of $\log s$. The unselected components after inflation cost at most
$s^{-1/48+o(1)}$. The projected old large-field site sum is bounded
by ${\cal W}(B_i)L_i^4T_eG(x_{i+1})$, also smaller than every inverse
power of $s$. It need not include $m$ a second time: the selected
component estimate already includes that factor.

The genuine additive multipliers have uninflated site bound
$C(42B_i+1)^4q_n$. Every fixed product of radius factors, support
conversion constants and $l$ is eventually at most $A_s^{1/4}$.
Hence their inflated site bound is at most $q_nA_s^{1/4}$.
For $n\le(\log\log s)^4$ and every fixed $r$,
\[
 q_nA_s^r\le s^{-1/400}.
\]
This follows by taking logarithms: the positive contribution is
$O((\log\log s)^7)$, whereas the negative term is
$-(\log s)/200$. The connected expansion can therefore be rooted
at an additive multiplier when one occurs; its site bound fits
$\omega_n=q_nA_s$. The other components fit
$\delta_n-\omega_n$, where $\delta_n=q_nA_s^3$, since
$s^{-1/48+o(1)}=o(q_n)$. This proves the inflated product bound and
its site summability on the genuine, not a substituted, density.
All finite stage factors in this paragraph are absorbed simultaneously
by its choice of a sufficiently large starting scale; no uniform supremum over infinitely many stages
is used.

It remains to reconstruct the coarse presentation. The regular
connected logarithm has anchored cost at most $16\delta_n$.
Keep the original coarse path action $A_i$. Its difference from
the base action extracted in the Gaussian integral consists of the
cubic, quartic and quadratic-loop terms restricted by their event indicators, with cost at most
\[
 \beta e^{7e}C_i\vartheta(x_i)^3+
 \beta e^{9e}C_i\vartheta(x_i)^4+
 e^{5e}\Gamma_i\vartheta(x_i)^2\le s^{-1/200}.
\]
The cover bound for a degree-$j$ path atom is
$1+2j+\operatorname{length}/B_i$; the path decay absorbs the last term.
Add this correction to the regular logarithm. Its total value cost
is at most $s^{-1/200}+16\delta_n\le q_{n+1}$.
The sign is fixed by
$e^{-\beta B_i^{\rm ref}}e^{-\ell_i}
=e^{-A_i}e^{-(\ell_i+\beta B_i^{\rm ref}-A_i)}$;
$B_i^{\rm ref}$ here is the base action, not the radius.

At the sliced-integrand level the event indicators retain the finite tag-count
factor; it cancels the denominator of the normalized tag average
when the slices are summed. Signed orientations retain their original
probability weights. Neither averaging introduces a second determinant
or a second extracted scalar.
The positive Gaussian prefactor determines a real scalar logarithm.
No branch of a complex logarithm is chosen. The remaining
large-field factors, including the regular logarithms with holes,
are grouped at their full support and exact old prescribed sets of large-field sites as in
\eqref{ch08:reassembly}. The raised hole correction is bounded by
$16l\delta_n e^{16\delta_n}$, using
$|e^{-c}-1|\le e^{|c|}|c|$.
This and the other component costs are $s^{-1/400+o(1)}$, within
the available $H(s)^{-1/8}$ reserve. The inflated connected bound
gives outgoing large-field bound at most $m^{-1}\le G(x_i)^2$,
hence at most $G(x_i)$. No residual path correction is left outside
the regular family. Averaging the regular pieces over the finite
signed motion group retains their anchored bound and invariant sum.
The exact factorization and reassembly thus give
$D_i=\mathcal B_{i,\beta}D_{i+1}$ on every coarse configuration,
including configurations with failed cutoff conditions. There is no division
by a large-field activity, which can vanish.

At each step choose this expansion before specifying the marked
family and retain it for subsequent use. Smaller support exponents
and smaller inflations are then estimates on the same coefficients.
Only the newly constructed lower entry is installed; the endpoint
and all already installed higher entries are unchanged. This proves
finite descending induction, including adjacent stage changes.
Finally positivity of $D_i$ follows from the positive endpoint and
probability kernels, not from positivity of individual activities.
Their mass-one identity and finite-volume integrability give
$\int D_i=Z_{N,p}$ by Fubini. The same equalities hold coefficientwise
for finite marked polynomials. This proves the asserted normalization
and the common-history statement.
\end{proof}


Finally let $C^{\rm ref}$ be a field-independent polynomial of prescribed
connected reference correlations with zero constant coefficient.
The corrected source is
\[
 Jq_N,\qquad J=\exp(-C^{\rm ref}),
\]
and its exact image under the block maps is $Jq_k$. Its coefficient of $z_\varnothing$ remains the
original density without insertions. One must start this comparison from the
full $q_k$, including all regular and large-field coefficients of source degree at least two.
Multiplying a freshly formed product of single insertions would
discard those coefficients and would not compare the original correlations.

\section{The normalized integral with several insertions}
\label{ch11:replicated}
Suppose the representation at the stopping scale, on $V$ sites, has $n$ copies,
weighted single-insertion bounds $x,y$, and higher regular and large-field
incidence at most one each. Assign each term with insertions to one of its
insertion points and each insertion-free term to one site in its support. The
coefficient norm is bounded by
\[
 e^{v_0+V(b_{\rm act}+\kappa)+n(x+y+2)}.
\]
Here $v_0$ is the constant term, $b_{\rm act}$ the absolute action cost per
site and $\kappa$ the insertion-free activity site sum. The estimate uses
$1+t\le e^t$ for the eligible activity product and retains every
coefficient of source degree at least two. The two unit bounds on the terms of source degree at least two come from
\eqref{ch11:higher-induction}.

Choose $0<\rho\le1$ below the small-field threshold. A maximal
packing in the ambient real matrix space of dimension $D_G$, with
Haar invariance, gives $\mu_G\{\chi(U)\le r\}\ge(r/4)^{D_G}$.
Restrict every forward link to normalized Hilbert--Schmidt distance from the identity at most $\rho/4$; the event has
probability at least $(\rho/16)^{4D_GV}$ and all plaquettes have
normalized Hilbert--Schmidt distance from the identity at most $\rho$. If the exact action there is bounded by
$V[b_{\rm act}+\beta b_{\rm loc}(\rho)]$, positivity gives
\[
 Z\ge e^{v_0-V[b_{\rm act}+\beta b_{\rm loc}(\rho)
                         +4D_G\log(16/\rho)]}.
\]
Signed regular costs are included in $b_{\rm act}$. Write $q$ for the full source density at the stopping level, as in
\eqref{ch11:source-density}, and $q^{\rm site}$ for the resulting
per-site cost. Dividing by the lower bound on $Z$ gives
\begin{align}
 \left|Z^{-1}\int[z_{\mathcal I}]q\,d\lambda\right|
 &\le\frac{e^{Vq^{\rm site}+n(x+y+2)}}{\prod_i\omega_i},\notag\\
 q^{\rm site}&=2b_{\rm act}+\kappa+\beta b_{\rm loc}(\rho)
                  +4D_G\log(16/\rho).
 \label{ch11:replicated-bound}
\end{align}
The constant term cancels. Proposition~\ref{ch11:mass} identifies the
coefficient integrals with the original ones. A reflected tiling at
scale $a_m$, stopped at $a_k$, has $(a_m/a_k)^4$ sites per cell.
Taking the number-of-cells root, with common weight $\omega_k$, gives
\begin{equation}
 {\cal R}_m(F)\le\omega_k^{-1}e^{(a_m/a_k)^4q_k^{\rm site}+x+y+2}.
 \label{ch11:root-bound}
\end{equation}
Here ${\cal R}_m$ denotes the chessboard seminorm of $F$ and
$q_k^{\rm site}$ the per-site cost at level $k$.
Multiplying an admissible residual weight by $c$ divides the bound
by $c$, uniformly in the period when the local bounds are uniform.


\subsection{An independent modulus of continuity in the insertion points}
\label{ch11:positional}

We use the usual reflection-positive transfer representation and
chessboard inequality for the Wilson measure
\cite{OsterwalderSeiler1978,FrohlichIsraelLiebSimon1978}.
These apply to the original positive Wilson law; no reflection
positivity of an individual effective activity is asserted.
The product inequality bounds insertions in distinct reflection cells
by their chessboard seminorms. Expanding a copied product of
$F+G$ over subsets of cells proves the triangle inequality for these
chessboard seminorms, including when one of them vanishes.

For correlation evaluation positions separated in every coordinate, the required
modulus follows independently of difference insertions. Represent a
cyclic correlation as a cyclic product of transfer operators
$\tr\prod_{i=1}^q B_iT^{H_i}/Z$, where $T$ is positive,
$B_i$ is a local slab kernel, the slab width is $w$, and the total
length is $M=qw+\sum_iH_i=2sw$.
The reflected-copy bound for a slab and its adjoint gives
$\|B_i\|_{2s}\le A Z^{1/(2s)}$.
The standard Schatten--H\"older inequality
\cite[Theorem~6.32]{TroppMatrixAnalysis2022} now bounds the product with
complex separations $h_i$ by
\[
 \left|\tr\prod_iB_iT^{h_i}\right|
 \le\prod_i\|B_i\|_{2s}\,
                 Z^{\sum_i\operatorname{Re}h_i/M}
 \le A^q Z ,
\]
provided their real parts are nonnegative and their sum equals the
original total free length. This uses the spectral identity
$\|T^h\|_{M/\operatorname{Re}h}=Z^{\operatorname{Re}h/M}$.

Moving one insertion by $l$ lattice units changes two separations to
$H_0-z,H_{q-1}+z$. If both original separations exceed $2R$,
the normalized product is analytic and bounded by $A^q$ for $|z|\le2R$.
Cauchy's estimate, integrated from $0$ to $l\le R$, gives a
difference at most $A^q l/R$. With $R=\rho/a_N$, this is
$A^q la_N/\rho$, a Lipschitz bound in continuum units, independent of the mesh.
The chessboard bound is chosen at one fixed interior cell scale below
the coordinate separation $\rho$.

The transfer kernels act on a compact continuous slice space.
The finite-matrix estimate passes to them by common tagged partitions:
replace each continuous kernel $K(x,y)$ by the matrix
$\sqrt{\mu(C_u)\mu(C_v)}K(\xi_u,\xi_v)$.
Positivity of $T$ is preserved. Every fixed cyclic trace converges
to its cyclic kernel integral by uniform continuity, including the
partition trace and the finitely many norms obtained from reflected copies of the slab.
Since $Z>0$, the integrated inequalities for the normalized products pass to the limit.
Thus the argument does not assume a finite-dimensional transfer space.
It gives the modulus at the separated correlation configurations needed
in Lemma~\ref{ch11:absorption}, without using the difference estimate
that lemma is intended to prove.

Recall that $\mathsf Z_{N,p,d}$ transfers a cutoff coefficient to
its fixed reference scale. A comparison coefficient is admissible
when its inverse coordinates under $\mathsf Z_{N,p,d}$, its \emph{cutoff-scale coordinates}, have a
uniform subpower bound. The graph norm still combines the regular
majorant with the large-field majorant multiplied by
$\kappa_m^{-1}=(\epsilon_m^{\rm for})^{-1/2}$.

\section{Backward comparison of a local trial}
\label{ch11:trial-descent}
The inhomogeneous graph estimate controls local trials as well.
Put $s_i=L_i^{1/8}$, $e_i=s_i^{-3}$ and $A_i=32(1+C_i)^2$.
Choose $100000K_i^6(1+C_{i+1})^2\le s_i$, so the raw error is
at most $e_i$. Subtract one normalized local insertion density from the trial at the ultraviolet cutoff,
choosing its reference coefficient so the residual coefficient at the coarsest scale
is zero. Invertibility permits this choice; exact linearity propagates
the same subtraction through every step.
Let $x_i$ be the residual coefficient and $c_i$ its full remainder norm bound, including the subtracted graph of normalized local insertions. With
$y_i=|S_{L_i}x_{i+1}|$, the raw inequalities give
\[
 y_i\le K_i\bigl(|x_i|+e_i(y_i+c_{i+1})\bigr),\qquad
 c_i\le K_iy_i+e_ic_{i+1},\qquad |x_{i+1}|\le y_i.
\]
If $A_i|x_i|\le c_i$, substitution gives
$y_i\le4K_ie_ic_{i+1}$. The ratio choice ensures
$4K_iA_{i+1}e_i\le1$ and $(4K_i^2+1)s_ie_i\le1$.
Hence $A_{i+1}|x_{i+1}|\le c_{i+1}$ and $s_ic_i\le c_{i+1}$.
The induction starts at the coarsest scale with zero coefficient.
If the trial at the ultraviolet cutoff has coefficient plus remainder cost at most $w$,
the bound on the graph of normalized local insertions gives $c_N\le w+C_N(w+|x_N|)$.
Absorb $|x_N|\le c_N/A_N$ to obtain $c_N\le2(1+C_N)w$.
Thus
\begin{equation}
 c_i+|x_i|\le4(1+C_N)w e^{-(\tau_N-\tau_i)/8}.
 \label{ch11:descent-error}
\end{equation}

For a trial $f=a_N^eF$ with fixed link support, extract its degree-$e$
coefficient $z$ and subtract the corresponding dimensionful section.
Split the difference with zero polynomial coefficient into its small-field part on the
union $D$ of the link supports and its large-field part. Attach the latter to
the original selected large-field factors; the terminal density is
unchanged. If $c_{N,p}$ bounds local derivatives and $b_{N,p}$
the global value, sufficient bounds are
\[
 q=K_ec_{N,p},\quad h=c_{N,p}+C_Nq,\quad
 r=b_{N,p}+2A_e\mathcal N_Nq,
\]
\[
 Q_N=1+q+e^{81\sigma_N}(h+2|D|r),\qquad
 E_{N,p}=4(1+C_N)Q_N+
       e^{81\sigma_N+162}r\zeta_{e,N}u_N^{-1}.
\]
Here $A_e$ bounds the global section, $\mathcal N_N$ is its finite
sum of inverse threshold powers, $\sigma_N$ the support exponent,
$u_N>0$ the strict margin for attachments at the ultraviolet cutoff, and $\zeta_{e,N}$ its
source weight. Eventually the fixed link support has connected covering number at most 81;
two extra support units contribute the factor $e^{162}$.
The first term absorbs the coefficient and small-field remainder; the
second absorbs the attachment of the large-field part with its original prescribed large-field set.
If $c_{N,p},b_{N,p},C_N,e^{81\sigma_N},\mathcal N_N,u_N^{-1}$
are subpower in $a_N^{-1}$, then so is $E_{N,p}$.

The backward comparison and the product-moment estimate now give, at every fixed
pair $(m,k)$ of cell scale and stopping scale, a reference label $v_{N,p}$ with cutoff-scale
coordinates bounded by $2E_{N,p}$ and
\begin{equation}
 \mathcal R_m\bigl(F-O_{N,p}(e,v_{N,p})\bigr)
       \le C_{e,m,k}E_{N,p}e^{-\tau_N/8}.
 \label{ch11:trial-error}
\end{equation}
Indeed the residual source weight is inversely proportional to its
bound in \eqref{ch11:descent-error}. All constants are independent
of period when the terminal norm bounds are. Thus the discrepancy for local trials tends to zero in the chessboard
seminorm. The same estimate applies to finite differences formed from
translates, using the derivative bounds on the translated link supports.
It bounds cutoff-scale coordinates, not the raw reference labels.

\section{Finite-dimensional detection before taking limits}
\label{ch11:detectors}

A useful part of the later difference argument is entirely finite
dimensional. We record it now, keeping its analytic hypotheses visible.
Let $V$ be a fixed finite-dimensional label space, $E_i\subset V$
the equation subspace at a finite cutoff, and $\Pi_i$ a bounded
projection with kernel $E_i$. Such projections need not converge.
Choose a fixed invariant Euclidean inner product and average an
orthogonal projection with this kernel over the finite symmetry group,
assuming $E_i$ is invariant. Each conjugate annihilates $E_i$ and differs
from the identity by a map into $E_i$. The average has these two
properties, which imply idempotence: if $v-\Pi_i v\in E_i$, then
$\Pi_i(v-\Pi_i v)=0$. Its norm remains bounded uniformly in $i$.

Let $L_i:V\to\CC^r$ be the finite family of correlation functionals used
to separate the quotient. Suppose
\begin{equation}
 |\Pi_i v|\le Z|L_i\Pi_i v|
 \label{ch11:detector-bound}
\end{equation}
with one constant $Z$, independent of cutoff and of the volume in continuum units.
This follows from a sufficiently accurate finite-cutoff comparison
with an injective free matrix of correlation pairings: if $|c|\le K|Gc|$ and
$\|Q-G\|\le(2K)^{-1}$, then
$|c|\le K|Qc|+\tfrac12|c|$, hence $|c|\le2K|Qc|$.
The separating family of free correlations comes from the finite Maxwell quotient; no
continuum rank assertion is needed for this argument.

\begin{lemma}[Absorption of a possibly large equation component]
\label{ch11:absorption}
Let $x_i\in V$ be encoded coefficients of a normalized lattice
difference and put $y_i=\Pi_i x_i$. Suppose $|L_i x_i|\le C$ and,
for
\[
 z_i=\frac{x_i-y_i}{1+|y_i|},
\]
the microscopic equation comparison gives
$|L_i z_i|\le(2Z)^{-1}$ eventually, uniformly in the volume in continuum units.
Then $|y_i|\le2ZC+1$.
\end{lemma}
\begin{proof}
Since $L_i y_i=L_i x_i-(1+|y_i|)L_i z_i$,
\eqref{ch11:detector-bound} gives
$|y_i|\le ZC+\tfrac12(1+|y_i|)$.
Move the last half-norm to the left.
\end{proof}

The normalized equation coefficient $z_i$ lies in $\ker\Pi_i$.
Applying an equation estimate to it still requires the estimate's
growth hypothesis. For instance, if $x_i$ has a subpower coefficient
bound, then $z_i$ has one: its first term is a bounded scalar multiple
of $x_i$, while $y_i/(1+|y_i|)$ is bounded by one.
Thus the argument does not presuppose a uniform bound for the raw
difference coefficient $x_i$.

There is also an exact way to test finite differences using bounded
separated observations. For a lattice spacing $a$, a coordinate unit vector $e$, and an integer
$h\ge1$,
\[
 \frac1h\sum_{j=0}^{h-1}
 \frac{F(x+(j+1)ae)-F(x+jae)}a
 =\frac{F(x+hae)-F(x)}{ha}.
\]
For a fixed physical distance $b>0$, taking $h=\lceil b/a\rceil$ bounds the right side by
$2\sup|F|/b$. To replace averaged correlation evaluation positions by the desired
positions, however, one needs a modulus of continuity in the insertion points established
independently of the normalized difference estimate being proved.
That modulus is an essential analytic input, not another consequence
of this telescoping identity. Likewise the microscopic equation
insertion must vanish against the distant spectator product
at finite cutoff, uniformly over the volumes in continuum units under
consideration. An identity obtained only after an ordered limit
would not justify the absorption lemma's uniform hypothesis.

\section{Collected estimates and local calibration}

The chapter now gives the quantitative single-insertion raw-map estimates
\eqref{ch11:graph-estimates}--\eqref{ch11:graph-difference},
including ball invariance and fixed-input continuity; the full
recurrence for terms of source degree at least two \eqref{ch11:higher-induction}; and the normalized
product-moment estimate \eqref{ch11:replicated-bound}.
The exact tree covariance \eqref{ch11:conditional-curvature} and the
local estimates in Section~\ref{ch11:actual-reader} identify the
joint Gaussian law used for the principal components.
The source weights and positive margins in the support exponent appearing in these
statements must be included in the selected stage requirements.

For a finite cutoff, composing the invertible single-insertion coefficient
maps from $N$ to the reference scale defines
$\mathsf Z_{N,p,d}$. It is real and symmetry-equivariant (the equivariance uses that of the insertion-free representations, which rests on Assumption~\ref{ch09:forest-assumption}, and the transport part of Hypothesis~(H) of Section~\ref{ch11:marked-step}; it is not proved here without them), and the exact
rule for the identity field gives its coefficient $a_{r(d)}^de_0$ at the reference scale.
The normalized local insertion is therefore
\[
 O_{N,p}(d,c;x)=a_N^{-d}
       P_d^{\rm loc}[S_{1/t_N}\mathsf Z_{N,p,d}^{-1}c](x),
\]
with the factor-degree scaling $S_t$ of the local section.
Uniform column estimates follow by multiplying the displayed
one-step inverse bounds; their subpower growth requires the
corresponding inequalities on the schedule.

There remains a distinct local calibration issue before claiming the
all-period finite-difference theorem. A translated finite difference
of degree $d$ has a protected local representative at degree $d+1$:
\[
 a_N^{d+1}\Delta_\alpha O(x)
 =a_N^dO(x+a_Ne_\alpha)-a_N^dO(x).
\]
This cancels the explicit $a_N^{-1}$, but moving the derivative-weight
origin contributes the factor $2^{\varpi_N q_N}$.
The lower bound on the starting scale of the stage must include
$2^{\varpi_Nq_N}\le\widehat\beta_N$ to keep that factor subpower.
The same requirement applies to repeated fixed finite differences.

Section~\ref{ch11:trial-descent} now gives the reference label $v_{N,p}$ of the trial, with
subpower cutoff-scale coordinates and a vanishing chessboard-seminorm
error, under its explicit terminal and stage bounds.
The microscopic equation section and its comparison with the dimensionful
section are given next. The argument with the separating family proves the finite-cutoff comparisons
needed for uniform differences.

% !TEX root = ../monograph.tex
\section{Calibrating the microscopic equations on the same history}
\label{ch11:equation-calibration}

Fix a degree $d$. The section $\mathcal V_d$ and automorphism $R_d$
are those of Theorem~\ref{ch02:equation-adapted}; neither is chosen
again when the cutoff or period changes. The physical section remains
$P_d^{\rm loc}$. We compare the two on the history already used to
define $\mathsf Z_{N,p,d}$.

We first specify an additional finite stage requirement.
At stage $j$, let $a_{dj}$ bound the derivatives of $\mathcal V_d$
requested by that stage, for every $d\le j$. These are finite constants
by the powered-curvature-factor estimate in Chapter~\ref{ch02:observables}.
Let $q_d$ bound the fixed union stencil and put
\[
 Q_{dj}=2+e^{81\sigma_j}(a_{dj}+C_j+2q_d).
\]
Here $\sigma_j$ is the support exponent and $C_j$ the column bound.
These quantities are fixed before the block length.
Let $\mathcal O_d\ge1$ bound both the finite symmetry actions on
coefficient $\ell^1$ norm and their averaging coefficients, and put
$\delta_d=(2\mathcal O_d^2)^{-1}$. Include the finite maximum
\begin{equation}
 L^{1/8}\ge
 \max_{d\le j}2(1+C_j)Q_{dj}\delta_d^{-1}
 \label{ch11:calibration-floor}
\end{equation}
in the stage-$j$ block floor. Include
$\widehat\beta_N\ge a_{d,J(N)}$ in the adjacent stage conditions on the starting scale. These requests are compatible with the selection order:
the section and its derivative constants are fixed first, then the
length, then the inverse coupling threshold. Raising a later threshold does
not change the section. Thus they can be imposed for all stages
before selecting the common history, with no dependence on $p$.

Define the normalized microscopic section at the origin by
\[
 O^{\mathcal V}_{N,p}(v)
   =a_N^{-d}\mathcal V_d[S_{1/t_N}v](0).
\]
The factor scaling $S_{1/t_N}$ commutes with $R_d$ only when
$R_d$ preserves factor count, which is not asserted here.
Accordingly its exact equation labels at this normalization are
\begin{equation}
 I_{N,d}^{\mathcal V}
    =S_{t_N}R_d\mathcal I_d.
 \label{ch11:scaled-equation-labels}
\end{equation}
Since $\mathcal I_d$ is homogeneous in factor count, one can equally
write $I_{N,d}^{\mathcal V}=
(S_{t_N}R_dS_{1/t_N})\mathcal I_d$.
This conjugation must be retained when using a cutoff-independent
unscaled equation section.

\begin{proposition}[One comparison map for all equation labels]
\label{ch11:equation-comparison}
Under \eqref{ch11:calibration-floor} and the terminal attachment
bounds of Section~\ref{ch11:trial-descent}, there is, for every
sufficiently large $N$ and every $p$, one real equivariant map
$P_{N,p,d}:V_d\to V_d$ such that
\[
 \|P_{N,p,d}-I\|_{1\to1}\le\tfrac12,\qquad
 \|P_{N,p,d}^{-1}\|_{1\to1}\le2.
\]
For each fixed stopping pair $(m,k)$,
\begin{equation}
 \mathcal R_m\!\left\{
 O^{\mathcal V}_{N,p}(v)
   -O_{N,p}(d,\mathsf Z_{N,p,d}P_{N,p,d}v;0)\right\}
 \le C_{d,m,k}\|v\|_1 E_N e^{-\tau_N/8},
 \label{ch11:equation-root-error}
\end{equation}
where $E_N$ is subpower in $a_N^{-1}$, uniformly in $p$.
The same map works for every spectator; it is constructed before
any spectator is specified.
\end{proposition}
\begin{proof}
For each basis vector, the exact terminal decomposition used in
Section~\ref{ch11:trial-descent} applies to the fixed section
$\mathcal V_d$. Its part with zero polynomial coefficient is localized on the union
stencil; the complementary part is attached to the original
large-field selection, with no change of density.
It gives coefficient $e_j$, ordinary norm bound
$Q_N=Q_{d,J(N)}$, and complete source norm bound
\[
 E_N=4(1+C_N)Q_N+
 e^{81\sigma_N+162}Z_N^{\rm val}\zeta_{d,N}u_N^{-1}.
\]
Here
$Z_N^{\rm val}=(A_{\mathcal V}+2A_{\rm loc})
 \sum_jt_N^{-s_j}$, with $A_{\mathcal V},A_{\rm loc}$
the global value norms of the two fixed sections.
The ordinary terminal norm bound uses a sharper attachment estimate
than the general trial bound. If $w_N^{\rm for}$ is the terminal
large-field bound and $D$ the union stencil, its contribution
to the remainder norm normalized by $\sqrt{w_N^{\rm for}}$ is
\[
 2|D|\sqrt{w_N^{\rm for}}\,
       e^{\sigma_N c_{B_N}(D)}Z_N^{\rm val}.
\]
Indeed attach the complementary local function to the original
nonempty prescribed large-field set and sum its possible distinguished
sites in $D$; the remaining large-field suppression bounds this
normalized mass. The same exact attachment retains reserved mass
$e^{(\sigma_N+2)c_{B_N}(D)}Z_N^{\rm val}$.
Impose the additional conditions on the starting scale
\[
 |D|w_N^{\rm for}\le1,\qquad c_{B_N}(D)\le81,\qquad
 \sqrt{w_N^{\rm for}}Z_N^{\rm val}\le1.
\]
Then $|D|\le q_d$ bounds the ordinary contribution by
$2q_de^{81\sigma_N}$, explaining exactly the definition of $Q_N$.
These are post-length threshold conditions. At fixed adjacent
stages the value norm is a fixed polynomial in the inverse
small-field threshold, whereas the large-field bound decays
faster than every such power. Their finite simultaneous threshold
therefore exists. The unweighted reserved mass is not discarded.
This last term is required: the ordinary derivative bound alone
does not bound the complementary large-field contribution.

Backward comparison gives a reference coefficient $v_j$ and the
chessboard-seminorm error in \eqref{ch11:trial-error}. Put
$u_j=\mathsf Z_{N,p,d}^{-1}v_j$.
There is a slightly stronger endpoint estimate than the displayed
bulk descent bound. At the last step its cone inequality gives
\[
 |e_j-u_j|\le
 4K_{N-1}L_{N-1}^{-3/8}c_N
 \le L_{N-1}^{-1/8}c_N
 \le 2(1+C_N)Q_N L_{N-1}^{-1/8}.
\]
The middle inequality follows from the previously imposed block
floor. Requirement \eqref{ch11:calibration-floor} bounds this by
$\delta_d$.

Let $Ue_j=u_j$. Average $U$ by conjugation with the finite geometric
symmetry group, and then with the real structure, obtaining $P$.
For the conjugation term use $\overline{U\overline v}$, so the result
is still complex linear. The norm of $P-I$ is at most
$\mathcal O_d^2\delta_d=1/2$, and the Neumann series gives the stated
inverse bound. Apply the same averaging to the field errors.
Linearity and the triangle inequality for the chessboard seminorm give \eqref{ch11:equation-root-error}, enlarging its fixed
constant to cover the finite orbit. This preserves one common
history, all source monomials, and one common map.

Finally $C_N$, $e^{81\sigma_N}$, $t_N^{-1}$,
$\zeta_{d,N}$, and $u_N^{-1}$ have the subpower bounds already
required for trial descent. The extra condition on the starting scale bounds
$a_{d,J(N)}$ by $\widehat\beta_N$, which is subpower.
Finite sums and products preserve this property, proving the claim
for $E_N$.
\end{proof}

The varying equation subspace in reference coordinates is therefore
\begin{equation}
 \mathsf E_{N,p,d}
   =\mathsf Z_{N,p,d}P_{N,p,d}
               S_{t_N}R_d\mathcal I_d.
 \label{ch11:actual-equation-subspace}
\end{equation}
For $v_N\in I_{N,d}^{\mathcal V}$ of subpower norm, the first term
in \eqref{ch11:equation-root-error} has exactly zero correlation
against every bounded spectator outside its microscopic stencil,
by \eqref{ch02:microscopic-null}. For fixed positive physical
separation, this exclusion holds at every sufficiently fine mesh,
uniformly in period. Holder's chessboard bound then shows that the
physical insertion with coefficient
$\mathsf Z_{N,p,d}P_{N,p,d}v_N$ has vanishing separated correlations
against any finite family of spectators with bounded chessboard seminorms. This proves the equation inclusion needed for the
argument using a finite separating family; equality with the full null space still
uses the free nondegeneracy and detection estimates.

The scaling in \eqref{ch11:actual-equation-subspace} also has a
useful check. Since $R_d-I$ never lowers factor count and $t_N\to0$,
the conjugate $S_{t_N}R_dS_{1/t_N}$ remains bounded.
Its strictly factor-raising terms tend to zero. Terms at equal
factor count can still raise scaling dimension and need not tend
to zero under this scaling alone. Thus the raw triangular
construction and the near-identity comparison $P$ do not, by
themselves, prove that this equation subspace approaches the
unmodified Maxwell ideal. No such identification is required here: use bounded projections
with kernel exactly \eqref{ch11:actual-equation-subspace} and establish
the bound from the separating family \eqref{ch11:detector-bound} on their complements.
The terminal map $\mathsf Z$ already includes the scaling-dimension
evolution; its physical convention has not been changed.


\section{The cutoff-dependent equation space and uniform finite differences}
\label{ch11:uniform-differences}

All eventual assertions in this section mean: there is $N_0$ such
that the assertion holds for every $N\ge N_0$ and every period $p$.
A bound after first fixing $p$ is insufficient.
Write $i=(N,p)$, $X_i=\mathsf Z_{N,p,e}$, and
\[
 K_i=X_iP_{N,p,e}S_{t_N}R_e\mathcal I_e.
\]
This is precisely the equation space
\eqref{ch11:actual-equation-subspace}. Every factor is invertible,
so $\dim K_i=\dim\mathcal I_e$ at each finite cutoff.
Choose the bounded equivariant projections $\Pi_i$ with kernel $K_i$
as in the construction of the separating family above, and fix
$A\ge\sup_i\|\Pi_i\|$. Their rank is
$r=\dim(V_e/\mathcal I_e)$, without any limiting rank assertion.

Call a coefficient family $z_i$ admissible here if
$\|X_i^{-1}z_i\|_1$ is bounded by a sequence subpower in $a_N^{-1}$,
independent of $p$. The inverse-column bound makes every bounded
coefficient family admissible. This class is closed under finite sums
and multiplication by scalar families of modulus at most one.

\begin{lemma}[Vanishing chessboard seminorms for the equation space]
\label{ch11:uniform-sweep}
If $z_i$ is admissible and $\Pi_i z_i=0$, its normalized local insertion
has chessboard seminorm tending to zero, uniformly in $p$, at every
sufficiently fine fixed cell scale. The assertion includes all its
finitely many signed and conjugated orientations.
\end{lemma}
\begin{proof}
Put $v_i=P_{N,p,e}^{-1}X_i^{-1}z_i$.
Its norm is subpower, and $v_i\in S_{t_N}R_e\mathcal I_e$.
The error in \eqref{ch11:equation-root-error} is therefore at most
$e^{-\tau_N/16}$ eventually, uniformly in $p$, after increasing
the fixed lower bound on $N$.

The microscopic comparison insertion itself has chessboard seminorm exactly
zero. Indeed, every other reflected copy lies outside its fixed
microscopic equation neighborhood once the mesh is sufficiently
fine relative to the fixed cell in continuum units. Apply
\eqref{ch02:microscopic-null} to one copy and the product of all
other copies. That product is bounded and measurable at each
finite cutoff, regardless of the size of its coefficients.
There is no limiting argument in this identity.
The smallest allowed period in continuum units is positive, so one lower bound
on $N$ ensures the required nonwrapping and separation for
every period. The triangle inequality for the chessboard seminorm proves the claim.
\end{proof}

In the following local comparison, $S_t$ multiplies each basis
coordinate by $t$ to its curvature-factor count, whereas $W_\rho$
uses scaling dimension. Neither is the coefficient map $\mathsf Z$.
The equation correction $R_d$ and its ideal $\mathcal I_d$ are
those of Theorem~\ref{ch02:equation-adapted}; the separate map
$P_{N,p,d}$ compares the two local sections.

\subsection{A common localization map for the separating family}
\label{ch11:common-extraction}

The map used to normalize free comparison columns cannot be chosen
independently of the polynomial that is integrated.
Here is the construction from the single-insertion representation.
Fix a polynomial degree $D\ge d$, to be increased below, before the
stage derivative order; require $D+1$ derivatives.
At the comparison scale $k$ put
\[
 s=\widehat\beta_k,\quad H=e^{(\log\log s)^2},\quad
 M=\lceil(\log s)^2\rceil,\quad B=M^2,\quad b=8B,
 \quad \gamma=\frac1{16(D+2)}.
\]
The original insertion norm uses radius $M$, whereas the restarted
representation uses radius $B$. These are different radii.

Let $A u$ be the actual local polynomial coefficient in dimensionful
normalization. The calibrated single-insertion bound controls the full
regular series by $C_k\|Au\|_1$.
Split it at original connected covering number $M$.
The part with covering number at most $M$ is contained in the
radius-$b$ patch: a connected chain of $M$ original-radius cubes
has diameter at most $4M^2+4M\le8M^2$.
The remaining part has norm at most
$\eta C_k\|Au\|_1$, with $\eta=e^{-eM}$, by retaining one
positive support exponent $e$.
Keep its original event indicators and all original large-field terms.
Only the near part is replaced below.

Use logarithmic link coordinates $U_e=\exp(\theta x_e)$,
$\theta=s^{-1/2}H$, in this patch. The Taylor polynomial through
ordinary degree $D$ of the near regular sum is well defined.
Indeed the series of derivatives is absolutely bounded by the
regular norm. Polarizing the directional derivatives bounds its
coefficient sum by $C_D(1+n_b)^D C_k\|Au\|_1$, where
$n_b=O(b^4)$ is the number of real link coordinates.
The Taylor remainder is bounded by
$C_k\|Au\|_1\|x\|^{D+1}$, up to a fixed constant.
The finite degree-$d$ localization by radial Taylor expansion of the near sum is its
prescribed local polynomial coordinate: all regular remainders have
zero extracted polynomial. Applying the extraction to a Taylor polynomial instead of the
holonomy curve itself changes this by at most
$C(1+b)^a\theta$ times its coefficient norm.
This follows from the local holonomy equation: the normalized
logarithm of a threshold holonomy and its linearized edge
integral differ by $\theta$ times a fixed-degree polynomial
in the link location. Only finitely many terms contribute
through degree $d$.

All constants just used are fixed before $s$.
A lower bound on the starting scale of the stage can impose
\[
 C_k\le H^\gamma,\qquad
 C_D(1+n_b)^D\le H^\gamma,\qquad
 C(1+b)^a\le H^\gamma.
\]
The second and third inequalities follow from
$\log b=O(\log\log s)=o(\log H)$; the first is a lower bound on the starting scale for
the already selected stage column bound.
The discarded-support factor $\eta$ and the nonlinear ray
error $\theta$ are smaller than every fixed inverse power of $H$.
Consequently the ordinary linear radial Taylor localization of the resulting
polynomial agrees with $Au$ up to relative error $H^{-1}$,
uniformly in $N,p,u$.

Convert links to linear sums of based plaquette curvatures in the
axial gauge on this same patch. Let $W_u$ be the resulting raw
polynomial in curvature divided by $g=\beta_{N,p,k}^{-1/2}$.
Each axial row has coefficient sum at most $4bg/\theta$.
The coupling band gives
$g/\theta\le\sqrt{8/7}\,H^{-1}$, so this row sum is eventually
at most one. Thus $\deg W_u\le D$ and
$\|W_u\|_{\rm coeff}\le H\|Au\|_1$, after the finite increases of the starting scale.
The construction is linear in $u$ because the Taylor polynomial
of the function itself is unique; it does not require a linear
choice of its separate activity representatives.

The raw polynomial also approximates the same normalized local insertion
on the restarted small-field region. Its threshold is
$\vartheta=s^{-1/2}e^{(\log\log s)^{3/2}}$.
Put $a=4b\vartheta$. The nonlinear axial link logarithm divided
by $\theta$ and the linear sum of plaquette curvatures divided
by $\theta$ differ by at most $C_Ga^2/\theta$, and their size is
at most $2a/\theta=H^{-1+o(1)}$.
Taylor's remainder and replacement of one monomial factor at
a time therefore give
\begin{equation}\label{ch11:raw-polynomial-error}
 |F_u^{\rm near}(U)-W_u(g^{-1}F(U))|
       \le H^{-D}\|Au\|_1.
\end{equation}
For example the Taylor term is at most
$4H^\gamma H^{(-1+\gamma)(D+1)}\|Au\|_1
 =4H^{-D-15/16}\|Au\|_1$.
The nonlinear coordinate error and discarded-support error
are smaller than any required inverse power.
For the solutions $u=u_{j,N,p}$ of the lift equations \eqref{ch11:simultaneous-detector-lifts} below, which satisfy
$\|Au\|\le H^{bd+c_d}$, choose $D>bd+c_d+v$ before the stage
orders to obtain error $H^{-v}$.

Let $K_k$ be the original same-insertion covariance of the
Gaussian comparison. Define the polynomial used in Wick calculations by
\[
 \mathscr P u=\exp\!\left(\tfrac12\Delta_{K_k}\right)W_u.
\]
This sign is deliberate: its evaluation with covariance whose
diagonal block is decreased by $K_k$ equals the Gaussian
evaluation of $W_u$ with the undecreased covariance.
The identity follows by multiplying the commuting finite
operators $\exp(\tfrac12\Delta_{K_k})$ and
$\exp(-\tfrac12\Delta_{K_k})$.
It remains valid for indefinite symmetric covariances.

Set $z=\theta/g\ge\sqrt{7/8}\,H$.
For $p_u(y)=W_u(zy)$,
\[
 (\exp(\tfrac12\Delta_{K_k})W_u)(zy)
   =\exp(\tfrac12\Delta_{z^{-2}K_k})p_u(y).
\]
The rescaled polynomial $p_u$ is the Taylor polynomial with
the factor $g/\theta$ in its axial substitution cancelled.
Its remaining axial rows have norm at most $4b$.
Since $(4b)^D\le H^\gamma$ eventually, its coefficient bound is
$4H^{3\gamma}\|Au\|_1\le H^\alpha\|Au\|_1$,
where $\alpha=4\gamma=1/(4(D+2))$.
The finitely many radial Taylor coefficient functionals likewise have coefficient
norm at most $H^\alpha$ after increasing the starting scale.
These bounds are needed before applying the Wick estimate.
The covariance itself has entry bound at most $H^{1/4}$ after
an increase of the starting scale, so the rescaled covariance entries are at most
$H^{-1}$ eventually.
For polynomials of degree at most $D$, the coefficient norm of
the difference between this Wick operator and the identity
is at most
\[
 H^{-1}w_D,\qquad
 w_D=\tfrac12D^2e^{D^2/2}.
\]
Indeed a contraction has coefficient norm at most $D^2/2$;
sum the finite exponential series, bounding its terms in absolute
value. This bound is independent of the number of ambient
coordinates because the polynomial coefficient norm already
sums their coefficients.

Let $\mathcal E$ send $v\in V_d$ to its polynomial, of weighted degree at most $d$, in the fixed radial jet variables, expanded in the fixed radial polynomial basis.
It is injective between finite-dimensional spaces, so one may choose $c_d\ge1$ with
$\|v\|_1\le c_d\|\mathcal Ev\|_{\rm coeff}$.
The extraction from $\mathscr Pu$ of the coefficients of low scaling dimension belongs to this basis:
the original covariance is a real spatial matrix tensored with
the identity on $\gLie$, so its contractions commute with the
adjoint action. Axial conversion and weighted truncation
also commute with that action. Define $Tu$ to be precisely
these basis coefficients, retaining the inverse coupling scaling above.
Using $\alpha=1/(4(D+2))$, the preceding coefficient bounds,
linear substitution, and the Wick estimate give
\[
 \|Au-Tu\|_1
 \le c_d\bigl(H^{-1}+D2^DH^{-1+\alpha(D+2)}
                  +w_DH^{-1+\alpha(D+1)}\bigr)\|Au\|_1
 \le\tfrac12\|Au\|_1
\]
eventually. Constants and conjugation are preserved at every
operation, hence $T1=1$ and $T\bar u=\overline{Tu}$.
This constructs the same real linear map for every column;
it is not an arbitrary replacement for $A$.
The inverse bound and simultaneous lift equations stated below
therefore apply to this common localization map onto the polynomial label space.

Finally let $\chi$ be the restarted small-field indicator on the
whole patch. Define $f_u^{\rm near,mask}$ to be only the near sum
with its original event indicators restored, including its constant term.
It equals $F_u^{\rm near}$ without its event indicator wherever $\chi=1$,
because the restarted threshold is smaller than the original one.
Then
\[
 f_u^{\rm near,mask}=\chi W_u+\chi(F_u^{\rm near}-W_u)
             +(f_u^{\rm near,mask}-\chi F_u^{\rm near})
\]
is exact. The second term has the bound
\eqref{ch11:raw-polynomial-error}; the third vanishes wherever
$\chi=1$ and is assigned to the original prescribed large-field set.
The distant regular sum and original large-field terms were
retained throughout. Substituting this identity into the finite
squarefree exponential preserves every higher source variable coefficient.
Thus the density, polynomial $\mathscr P$, common
localization map $T$, and normalized lifts all refer to the same
terminal insertion.


The common localization map may also be chosen equivariant under the
finite signed coordinate group. Average the raw polynomial
construction on domain and range before defining $T$.
The insertion, covariance and radial Taylor localization have the
corresponding covariance identities, so this preserves both
the exact density comparison and its extraction of the coefficients of low scaling dimension.
All errors are enlarged by at most a fixed finite symmetry
norm. Since their relative bounds tend to zero, choose the
threshold to make the averaged error at most one half.
Average with conjugation at the same time. This uses one
polynomial construction for all transformed columns, rather
than symmetrizing independently chosen lifted vectors.

At a fixed comparison scale $k$ the one-step coefficient bounds in the direction of the transfer, applied successively from the reference scale to $k$, give
$\|u\|_1\le D_k\|Au\|_1$, uniformly in $N,p$.
The relative half-error therefore implies
$\|u\|_1\le2D_k\|Tu\|_1$, so $T$ is injective and hence
invertible. Let $v_1,\ldots,v_r$ be real representatives of the
free separating family. Write $W_\rho$ for scaling-dimension
dilation and $S_t$ for factor-count dilation. With
$\rho=H^b$ solve all the equations
\begin{equation}\label{ch11:simultaneous-detector-lifts}
 S_{g/\theta}T u_{j,N,p}=W_\rho v_j.
\end{equation}
They have unique real solutions using the same $T$.
The coupling band gives $\theta/g\le2H$, and therefore
\[
 \|u_{j,N,p}\|_1
 \le2D_k\|S_{2H}W_\rho v_j\|_1.
\]
Their sum is a finite bound $B_k$ independent of cutoff and
period. No bound uniform in the comparison scale is needed.
Moreover $\|Au_j\|\le2\|Tu_j\|$ is bounded by a fixed
power $H^{bd+c_d}$, where the additional exponent $c_d$
depends on factor counts and the finite target list, not on
$b$. This is the growth factor used when choosing the
Taylor order and the geometric magnification.

Here $c_B(X)$ is the connected covering number at radius $B$.
The support exponent multiplies $c_B(X)$ in an exponential norm,
whereas a source weight is raised to the number of distinct
insertion indices. Changing one of these weights does not change
the coefficient or its original prescribed large-field set.

\subsection{Absorbing the larger source weight}
\label{ch11:weight-restoration}

A bound by one at the original source weight is not enough for the
comparison of separating correlations. The following argument absorbs the necessary change
of weight uniformly in the number of reflected cells. Write $n\le2$
for the number of insertions in each cell, fix $U\ge1/16$, and put
\[
 \omega_-=H^{1/16-U},\qquad \omega=H^{1/16}.
\]
For an activity with insertions, with support $X$ and source variable set $I$, its
weighted absolute bound contains $\omega_-^{|I|}e^{e_*c_B(X)}$.
Choose $e_*=e+2$ before constructing the expansion with insertions.
The same expansion, estimated at $e_*$, is used below; there is
no change of coefficients when its support exponent is lowered.

Each prescribed support is contained in the fixed enlargement of a
radius-$AB$ ball around its cell's finitely many insertion points, with $A$
chosen before the inverse coupling. A connected grid of radius-$B$ cubes
covers it with at most $C(1+A)^4$ cubes. Thus its support bound
at exponent $f$ is bounded by
$R_f=\exp\bigl(fC(1+A)^4\bigr)$, independent of $B$ and the period.
At successive stages these are finite pre-threshold constants.
The stage choices of sufficiently large starting scales impose a common bound $H^{\epsilon_j}$
on $R_{e+2}$, $C_*$, $L^4T_{e+2}$, and the fixed incidence
constants, with $\epsilon_j\downarrow0$.
Consequently every finite product of these constants is
$H^{o(1)}$ along the selected comparison scales.
This uniform lower bound on the starting scale, rather than treating a changing stage
constant as an absolute constant, justifies the exponent
comparisons below.
This bounded-cover assertion is used at the middle scales.
It is not an assertion that the separated opening sites
at distance $H^b$ fit in one radius-$AB$ ball.
The supports of their terms of source degree at least two are instead estimated by
\eqref{ch11:opening-higher-weight}.
At an earlier step of the \emph{upper interval}, from the comparison scale to the middle scale, a nonzero connected principal component
has radius at most $(6n+3)B$ around each of its own insertion points;
a term of source degree two with more widely separated insertion points is absent.
Only these nonzero components enter the principal support bound.

Let $Y_c$ contain all prescribed supports in cell $c$. Suppose
these sets have mutual distance at least $D'$, and every activity
support contains the sites of all its insertion points. Let $t(I)$ count
the cells represented in $I$, and let $t(X)$ count the sets $Y_c$
met by $X$. Then $|I|\le nt(I)\le nt(X)$.
For $t(X)\ge2$, a connected $B$-cover gives
\begin{equation}\label{ch11:dispersal}
 c_B(X)\ge\lambda t(X),\qquad
 \lambda=\frac{D'}{3(2B+1)},
\end{equation}
once $D'/3>2B+1$.
To see this, choose one point of $X\cap Y_c$ per represented cell,
and follow a path in its connected cover until distance
$\lfloor D'/3\rfloor$ from that point. A cube covers at most
$2B+1$ of the successive distance values. The resulting initial
path segments belong to disjoint distance balls, and a cube
cannot meet two of them. Summing proves the inequality.
Here connected cover means the connected cover used in the
activity norm; one cannot replace it by an arbitrary disconnected
cover of the insertion points.

\begin{lemma}[Regular remainders at the larger source weight]\label{ch11:restore-regular}
Suppose the regular remainder has incidence bound at most
$H^{\kappa-q}$ at $(\omega_-,e+2)$, with support containing its
insertion points, and $\lambda/\log H\to\infty$.
If $q\ge\kappa+Un+2$, its incidence bound at $(\omega,e)$ is at most
$H^{-1}$ for all sufficiently large $H$.
\end{lemma}
\begin{proof}
Separate activities meeting one cell from those meeting at least
two. The first class gains at most $H^{Un}$ on changing weight.
For the second class retain one additional unit of the support exponent and use
\[
 H^{U|I|}e^{-c_B(X)}
 \le H^{Unt(I)}e^{-\lambda t(I)}\le1
\]
eventually. Thus the new bound is at most
$(H^{Un}+1)H^{\kappa-q}\le2H^{-2}\le H^{-1}$.
The estimate is independent of the total number of cells.
\end{proof}

The support containment is preserved by the algebra.
Initially an insertion contains its site. A coarse activity
contains the projected union of the fine supports, and its source variable
set is their union. Prescribed support enlargement, event indicator splitting
and the connected regrouping all retain those projected points.
This induction applies to both regular and large-field activities.
For a window whose cell side at its lower comparison scale is at least
$(\log s)^{20}$ times the middle lattice spacing, whereas
$B\le4(\log s)^4$, the cell neighborhoods have
$D'\ge(\log s)^{20}/2$ after their fixed-radius enlargements.
With $H=\exp((\log\log s)^2)$ this proves
$\lambda/\log H\to\infty$, uniformly in period.
The condition concerns the distance between replica cells, not
between the two insertions within one cell.

\begin{lemma}[Large-field activities at the larger source weight]
\label{ch11:restore-forced}
Suppose the old large-field site bound is $\eta_-$ at
$(\omega_-,e+2)$, where $\eta_-$ decays faster than every inverse
power of $H$. Suppose each cell has at most $2^n$ additional
\emph{failed-cutoff factors} (the terms $-(1-\chi)Rz_I$ of the factors $1-(1-\chi)Rz_I$ in the proof below, which vanish where the cutoff $\chi$ equals one), supported in $Y_c$, whose coefficient
bounds before source variable weighting are at most $C_*=H^{o(1)}$ and whose
support bounds at exponent $e+1$ are at most $R_{e+1}=H^{o(1)}$.
Every nonzero such factor must meet an old large-field activity.
Then their exact connected regrouping has site bound at most
\[
 4H^{Un+2^n}\eta_-
\]
at $(\omega,e)$, for all sufficiently large $H$.
\end{lemma}
\begin{proof}
A failed-cutoff factor of nonconstant source monomial $I$ occurs in the exact
identity
\[
 \exp_*(\chi W-R)
   =\exp_*(\chi(W-R))\,\exp_*(-(1-\chi)R).
\]
Its last exponent is square zero for each source monomial.
Expand it as a product of factors $1-(1-\chi)Rz_I$.
Whenever such a factor is nonzero, its support contains a large-field indicator
covered by an original prescribed large-field set. Thus every surviving
connected component contains an old large-field activity.
This is a pointwise identity of the full density with insertions.

Put $F=\log H$. A new failed-cutoff factor has bound
$p_f\le H^{1/8}C_*R_{e+1}$ and assign it auxiliary bound
$w_f=e^{-F}p_f$. Assign an old activity $o$ the auxiliary bound
\[
 w_o=e^g e^{2^nFt(X_o)}p_o,
\]
where $p_o$ is its bound at $(\omega,e+1)$.
The original bound $p_o^-$ at $(\omega_-,e+2)$ gives
\[
 w_o\le e^gH^{(Un+2^n)t(X_o)}
                     e^{-c_B(X_o)}p_o^-.
\]
For activities meeting at least two cell neighborhoods,
\eqref{ch11:dispersal} absorbs the power of $H$.
For the remaining activities its factor is at most $H^{Un+2^n}$.
The old auxiliary site sum is therefore at most
$e^gH^{Un+2^n}\eta_-$.

Set $\delta=H^{-1/2}$ and choose
$e^g=\delta/(2H^{Un+2^n}\eta_-)\ge1$.
The new factors have auxiliary site sum at most
$2^nC_*R_{e+1}H^{-7/8}$, so the total auxiliary sum
$\Gamma$ is at most $\delta$ eventually.
Assign each failed-cutoff factor to an old activity it meets.
At most $2^nt(X_o)$ factors are assigned to $o$.
In each nonzero connected component, their discarded
$e^F$ factors are thus bounded by the $e^{2^nFt(X_o)}$
factors, leaving at least one factor $e^{-g}$.
The connected-set estimate of Chapter~\ref{ch08:chapter}
gives the site sum at exponent $e$ at most
\[
 e^{-g}(e^\Gamma-1)
 \le2\delta e^{-g}=4H^{Un+2^n}\eta_-.
\]
Its counting condition $(2B+1)^4\delta\le1$ follows from
$B=O((\log s)^4)$ and the definition of $H$.
No independence of large-field indicators or activities is used.
\end{proof}

These lemmas explain the required stronger support estimate.
For the same expansion on the upper interval let its regular remainder satisfy
\[
 z_0\le H^{-q},\qquad z_{j+1}\le(\alpha+2)z_j,\qquad
 (\alpha+2)^{k-\ell}\le H^\kappa
\]
at exponent $e+2$ and weight $\omega_-$.
This implies the premise of Lemma~\ref{ch11:restore-regular}.
The constants at exponent $e+2$ must be fixed before $\alpha$
and the block ratio. Applying an estimate proved only at exponent
$e$ a second time does not establish this premise.

If the corrected principal coefficients at the middle scale are at
most $H^{-31/32}$, their contribution to the incidence bound at the larger source
weight is at most $2^nR_eH^{-27/32}$.
Together with the two preceding lemmas this gives regular
bound
\[
 H^{-1}+2^nR_eH^{-27/32}\le H^{-3/4}
\]
and large-field bound at most $H^{-q}$.
The reference connected polynomial is subtracted once, on the
same sample and with the same subset moments.
Its large-field part is precisely the part handled above.

On the lower interval suppose the step with insertions has the
quantitative bound
\begin{equation}\label{ch11:lower-weighted-step}
 z_{j+1}\le\alpha z_j+H^{1/2}z_j^2
\end{equation}
at the larger source weight, and its large-field part stays in the
allowed interval. If
$(\alpha+1)^{\ell-w}\le H^{1/4}$, induction gives
\[
 z_j\le(\alpha+1)^jH^{-3/4}\le H^{-1/2}.
\]
Thus every step remains in the domain
$H^{1/2}z_j\le1$ used in the induction.
At the output the source variable incidence and large-field site bounds
are at most one. On a torus with $q_0$ sites and $nq_0$ insertions,
the scalar majorant of the exponential with insertions gives
\begin{equation}\label{ch11:weighted-output-envelope}
 \omega^{q_0|I|}|[z_I^{\rm all}]\rho_{\rm marked}(U)|
 \le\exp\bigl(\text{unmarked exponent}(U)+q_0(1+n)\bigr).
\end{equation}
Every cell carries the source monomial $I$ in the coefficient on the left.
The integral without insertions is still the original partition function
by \eqref{ch11:source-density}.
This gives the numerator needed in the finite separating family
comparison, once the stronger upper recurrence and
\eqref{ch11:lower-weighted-step} have been proved for the combined estimate with insertions. The bound by one of \eqref{ch11:higher-induction} for terms of source degree at least two has neither
the source weight nor the small initial remainder needed here.

\paragraph{The two weighted step estimates.}
The inflations required here depend on the inverse coupling, so the
fixed-inflation statement alone is insufficient.
Split the decay exponent $Kh(x)^2r$ of a selected component with $r$ selected cutoff conditions, given by
Chapter~\ref{ch06:large-fields} (Section~\ref{ch11:restart-construction}), into fractions kept aside and choose
the reported bound $G(x)=\exp\{-a_*h(x)^2\}$ with $4a_*\le K/2$ so that multiplication by $G(x)^{-2}$
leaves a fixed positive fraction. This choice is made before the
thresholds. The amplitudes without prescribed large-field sites have a negative inverse coupling power
times subpower factors; multiplying by $H^c$ for fixed $c$ retains
that power, since $\log H=o(\log x)$. These margins establish the common
product bound \eqref{ch11:integrated-component} for the following
inflations, including their product over distinct prescribed large-field sets.

Use $e_*=e+2$, put $T=T_{e_*}$ and choose $\alpha\ge8T$.
For $H^{-q}\le z\le H^{-4}$ and principal incidence at most $H^{-4}$,
take
\[
 l=\max\left(1,\frac{H^{-3}}{\alpha z+H^{-q}}\right),
 \qquad m=G(x)^{-2}.
\]
Then $l\le\max(1,H^{q-3})$ and $l/x\le1$ eventually.
Thus the actual probability of a large fluctuation-field event absorbs the complementary
principal cutoff. The insertion incidence satisfies
\[
 a=lTz+2TH^{-4}\le\tfrac12H^{-3}.
\]
If $l>1$, use $lTz\le(T/\alpha)H^{-3}$; if $l=1$, use
$z\le H^{-4}$. Since $\mathcal W(B)=H^{o(1)}$ and $G(y)$
decreases faster than every inverse power of $H$, choosing
$\delta=H^{-1/8}/16$ gives $\Sigma\le H^{-1/8}/4$ and
$TG(y)+a\le H^{-3}$.
The connected bound and division by $l,m$ give
\[
 z_{\rm nonprincipal}\le H^{-3}/l
       \le\alpha z+H^{-q}\le(\alpha+1)z,\qquad
 z_{\rm principal}\le H^{-3},\qquad
 z_{\rm forced}\le G(x)^2.
\]
When $l=1$, its definition implies
$H^{-3}\le\alpha z+H^{-q}$, so the first estimate still holds.
Reassigning a principal term to its prescribed enlarged support
adds at most $2^nR_{e_*}y$ if its coefficient difference is $y$.
The joint Gaussian estimate at the requested power gives
$2^nR_{e_*}y\le H^{-q}$.
Its large-field part retains the original
prescribed large-field set. With $e^{-g_i}=G(x)$ the bound
\[
 e^{g_i}G(x)^2+nH^{-2}R_{e_*+1}\le H^{-1}
\]
gives outgoing large-field bound
$G(x)(e^{H^{-1}}-1)\le G(x)$.
This proves the upper recurrence
$z_{\rm new}\le(\alpha+2)z$ at the stronger exponent.

Below the middle scale there are no principal factors.
Use exponent $e$, $l=1$, $m=G(y)^{-1}$ and incoming large-field
bound at most $z\,a_{\rm low}(y)$, where
$0\le a_{\rm low}(y)\le1$. The outgoing incidence is at most
\[
 T_ez(1+a_{\rm low}(y))
 \le2T_ez\le\tfrac12\bigl(\alpha z+H^{1/2}z^2\bigr).
\]
The site conditions hold for $H^{-q}\le z\le H^{-1/2}$
by the same estimates, and the large-field output is at most $G(y)$.
This proves \eqref{ch11:lower-weighted-step}.
For the polynomial attenuation used here, the exponential
large-field decay kept aside gives $G(y)\le H^{-q}a_{\rm low}(y)$. Since
$z\ge H^{-q}$, the outgoing large-field bound is at most
$z\,a_{\rm low}(y)$, as required for the next step.

This calculation requires the decay kept aside and the inflated product bound
\eqref{ch11:integrated-component} for the factors without insertions of the restarted representation.
It does not transfer a bound from a different Gaussian law or
change an incoming bound after the history is fixed.
The restarted representation keeps the original path and clipping coefficients and Gaussian law (Section~\ref{ch11:restart-construction}).

The parameter $H=e^{(\log\log s)^2}$ measures the slowly growing
comparison scale, not a covariance or a Hilbert space.
The Gaussian fluctuation scale at the inverse coupling $\beta$ is $g=\beta^{-1/2}$.
The density $q_k$ with insertions remains the full transferred polynomial of
\eqref{ch11:source-density}, including its higher coefficients;
it is unrelated to the fixed exponent $q$ of the initial remainder bound.

\subsection{The covariance on a comparison window}
\label{ch11:window-covariance}

We record the covariance estimates used in the comparison of separating correlations,
including their finite-period errors. Fix enough smoothing orders for
the finite number of differences below. Let $H(p)$ denote the original
two-form metric symbol of Chapter~\ref{ch05:gaussian}, with
$H(0)=I$, a positive lower bound, and uniform exponential coefficient
bounds. Put $q_\mu(p)=e^{ip_\mu}-1$, let $X$ be exterior
multiplication by $q$ from one-forms, and let $U$ be exterior
multiplication from scalars. The curvature covariance is
\[
 P(p)=X(p)\bigl(X(p)^*H(p)X(p)+U(p)U(p)^*\bigr)^{-1}X(p)^*,
 \qquad P(0)=0.
\]
Indeed the completed inverse differs from the original potential
inverse only in gradient directions, annihilated by $X$.
The exterior identity gives coercivity at least
$c|q|^2$. Taylor expansion at zero, with differentiated bounds, yields
\[
 D^jP(p)=O(|p|^{-j}),\qquad
 D^j(P(p)-P_0(p))=O(|p|^{1-j}),
 \quad P_0=X_0|p|^{-2}X_0^*,\quad X_0=ip\wedge.
\]
All constants are uniform in the kernel history within the
already fixed quadratic bounds.

Use a smooth compact frequency partition whose $2\pi\ZZ^4$
translates sum to one, equal to one near zero. This leaves the
finite Fourier sum unchanged, including its zero-frequency value.
For difference orders $j,h$, subtract the Gaussian-regularized
continuum multiplier. The remainder has order $j+h+1$ at zero
and rapidly decreasing derivatives away from zero.
Splitting its Fourier integral at radius $|x|^{-1}$ and
integrating by parts $j+h+6$ times outside that ball gives
$C|x|^{-j-h-5}$. The inner integral has the same bound.
The finite Fourier sum differs from its integral by at most
$C(1+|x|)/N_{\rm per}$: on frequency cubes of side
$N_{\rm per}^{-1}$ use the derivative bound $C/|p|$,
which is integrable in four dimensions, and the phase difference.
This also handles the frequency cube containing $p=0$.
Gaussian regularization of the continuum potential has kernel
$(4\pi^2|x|^2)^{-1}(1-e^{-|x|^2/4})$ in radian convention.
Differentiating it therefore changes the separated Maxwell
covariance only by an exponentially decreasing term.
Consequently, for $|x|\ge1$,
\begin{equation}\label{ch11:finite-covariance}
 |\Delta^j\Delta^h\kappa_k(x)-\partial^j\partial^h C_{\rm Max}(x)|
 \le C|x|^{-j-h-5}+C(1+|x|)/N_{\rm per}.
\end{equation}
The lattice-side lower bound for the smallest period in continuum units
makes the second term negligible on every fixed-power-$H$ window,
uniformly in all periods.

For the installed curvature response $\Psi_i$, one also
needs to move differences through the curvature response, not merely bound
its differentiated row. The box-multiplier identities in the alias
formula give
\[
 \Delta_w\Psi_i=Q_{i,w}\Delta_w,\qquad
 \sup_x\sum_y|Q_{i,w}(x,y)|\le R L_i^{-|w|-2}.
\]
On the principal alias divide by the coarse difference using
$(e^{ip/L}-1)/(e^{ip}-1)=
[L\,\operatorname{avg}_L(p/L)]^{-1}$.
On other aliases consume one box factor for each difference;
their geometric-series identity has the same coarse factor.
Sufficient smoothing order leaves the remaining powers
nonnegative. The quotient is regular also on coordinate
hyperplanes. Its exponential coefficient bound follows at a
narrower decay width from
$(e^{inp}-1)/(e^{ip}-1)=\sum_{a=0}^{n-1}e^{iap}$
for $n>0$, with the analogous negative sum for $n<0$.
The extra factor $|n|$ is absorbed by that strict margin in the decay width.
This proves the row bound with $R$ fixed before $L$.
The curvature response preserves closed cochains, including harmonic ones,
by the exact exterior-transport identities of Chapter~\ref{ch05:gaussian}.

Let $\mathcal L=\Psi_{k-1}\cdots\Psi_\ell$,
$m=k-\ell$ and $S=a_\ell/a_k$.
A sign-valued test function on a specified row, without requiring that function to be closed, gives
\[
 S^{j+2}\sum_y|\Delta_w\mathcal L(x,y)|
       \le2^d R^m\quad(j=|w|\le d-1).
\]
Choose $m=\lceil(b+2)\log H/\log L\rceil$ with $b\ge6$,
and choose $L$ after $R$ so that
$4(b+2)\log R\le(\lambda/4)\log L$.
Then $S\ge H^{b+2}$ and the right side is at most
$P=H^{\lambda/4}$ after an increase of the starting scale.
The middle-scale field gives such a closed background field,
rather than an independently chosen one. In axial gauge on its
radius-$B$ patch take $g^{-1}$ times the Lie-algebra curvature
coordinate of each link, extend this one-cochain by zero, and
take its discrete exterior derivative $u$. Then $du=0$ exactly,
including across the artificial patch boundary, and
$\|u\|_\infty\le16B\vartheta/g=H^{o(1)}\le P$.
The truncation at the patch boundary is included in the
comparison error. Retain only response paths lying in the protected
patch and put the omitted paths in the regular remainder.
The exponential response bound makes their absolute contribution
at most $C B^a e^{-cB}g^{-1}$, with the same bound for every fixed
requested difference. The clipped field outside the patch is
bounded before multiplication by $g^{-1}$.
Since $B$ grows as $(\log s)^4$, this error is smaller than every
fixed inverse power of $H$, even after $g^{-1}=O(\sqrt s)$.
Thus zero extension is not asserted to preserve the full
conditional mean exactly.
On the interior stencil its difference from the nonlinear normalized
plaquette curvature is $O(B^2\vartheta^2/g)$, smaller than every
fixed inverse power of $H$. This follows by subtracting the
linear sum from the four-link product in that same gauge.
For this closed middle-scale background field $u$,
the mean $\mu=\mathcal Lu$ and the conjugated coarse covariance
$B_0=\mathcal L\kappa_\ell\mathcal L^T$ satisfy
\[
 S^{j+2}|\Delta^j\mu|\le P^2,\qquad
 S^{j+h+4}|\Delta^j\Delta^hB_0|\le P^3.
\]
The exact conditional covariance identity
\eqref{ch11:conditional-curvature} telescopes through the window.
After the same-insertion normal ordering used for
$\mathscr P$, its covariance is
\[
 \Gamma_{ij}(z,w)=\mathbf1_{i\ne j}
       \kappa_k(z_i+z,z_j+w)-B_0(z_i+z,z_j+w).
\]
In particular its diagonal is $-B_0$, not zero.
It need not be positive; only the polynomial Wick evaluation
uses this covariance.

Take the measurement magnification $\rho=H^b$.
The preceding bounds, \eqref{ch11:finite-covariance}, and
$\rho/S\le H^{-2}$ imply weighted mean and covariance errors
at the retained jet coordinates at most $H^{-3}$.
The highest difference and remainder coordinates have bound $P$.
Here the continuum approximation error contributes
$O(\rho^{-1})=O(H^{-b})$, the mean contributes
$H^{-4}P^2$, and the conjugated covariance contributes
$H^{-4}P^3$. The finite-period term is smaller than these
after the common lower bound on $N$.

To turn these into one finite coordinate system, choose dual
finite-difference functionals $\ell_\alpha$ on the polynomial
closed curvature cochains $t_\alpha$, with
$\ell_\alpha(t_\beta)=\delta_{\alpha\beta}$.
They exist by downward induction from the highest nonzero homogeneous
curvature, exactly as in the radial-jet construction.
Retain coordinates $\rho^{g_\alpha}\ell_\alpha f$ and,
at every stencil point $z$, the residual coordinate
$\rho^{d+1}(f-\sum_\alpha\ell_\alpha(f)t_\alpha)(z)$.
This gives an exact reconstruction of every stencil value.
For a closed cochain, repeated finite-difference summation along
coordinate paths bounds its residual by the top difference times
a fixed power of the stencil radius.
These factors are $H^{o(1)}$ and are absorbed by $H^{\lambda/4}$.
The common enlarged mean and covariance are therefore bounded by
$H^\lambda$, and their retained jet block differs from the free
one by at most $H^{-2}$. The free comparison has zero residual
coordinates. All subsets use this one enlarged system.

\subsection{From the common localization map to the corrected pair bound}
\label{ch11:corrected-pair-comparison}

Solve the simultaneous lift equations using the map $T$
constructed in Section~\ref{ch11:common-extraction}.
For the finitely many free polynomials, their coefficient sizes
are fixed. The solved lift has a bound $H^{bd+c_d}$ in normalized local insertion
coordinates. Choose $b$ larger than $c_d+3$, then choose $D$
large enough for \eqref{ch11:raw-polynomial-error}; all these
choices precede the derivative orders and final scale thresholds.
After the common finite-coordinate substitution, the retained
weighted polynomial is exactly the intended free polynomial.
Every remaining monomial has scaling dimension at least $d+1$, so its
coefficient sum is at most $H^{-2}$ by the factor
$\rho^{-(d+1)}$ and this choice of $b$.

For completeness, if a polynomial has coefficient norm $A$ and
degree at most $D_0$, Wick recursion bounds its expectation by
$A[2B(D_0+1)]^{D_0}$ when all means and covariance entries are
bounded by $B\ge1$. Changing these entries by at most $\epsilon$
changes it by at most the same expression times $\epsilon$.
This is proved by choosing the first variable of a monomial
and summing its mean term and at most $D_0-1$ pair terms;
subtract products one factor at a time.
It is independent of the number of coordinates.

Set $\lambda=1/[4(2D+D+1)]$ and $B=H^\lambda$.
The common coordinate estimate and the coefficient tail just
proved give, for every subset of the two insertions, moment
error at most $C H^{-2}H^{2D\lambda}\le H^{-3/2}$.
The empty moment is exactly one.
The nonlinear-to-polynomial error has the same bound:
use \eqref{ch11:raw-polynomial-error}, the
Gaussian estimate under the joint law with a smooth cutoff, and telescope the finitely many polynomial
transfer steps. If a step has coefficient norm at most
$A_{2D}$ and difference at most $H^{-v}$, the product error
over $m$ steps is at most $mH^{-v}A_{2D}^m$.
Choose the block length after $A_{2D}$ and then $v$ so this is
at most $H^{-3/2}$. The local-coordinate error is a negative
inverse coupling power times subpower factors and hence is below every
fixed $H^{-v}$ on the comparison band.
The comparison uses the same mean, covariance and common
cutoff factors at every subset; independently chosen Gaussian
samples would not justify the product estimate.

For two source variables the connected coefficients are simply the two
means and the covariance of the polynomial values. Subtraction
of these expressions gives a fixed polynomial bound in the
subset moments times their errors. Their free values remain
bounded on the chosen separated configurations.
Thus the corrected principal coefficients at the middle scale
are at most $H^{-31/32}$, with room to spare.
Here is the explicit estimate for the terms of source degree at least two already present
at the opening scale. Write $C=\lceil(\log s)^2\rceil$, $B=C^2$,
let $\theta_0>0$ be their original support exponent, and let
$r_0=s^{-a/2}$ be their original source variable radius. Their normalized
rooted bound is at most one by \eqref{ch11:higher-induction}.
The new radius is $r=H^f\le1$, where $f=1/16-U\le0$.
For a support containing $m\ge2$ distinct insertion points, the
chosen configurations have pairwise distance at least $C^2$.
The separated-site bound and cover enlargement give
\[
 Cm\le64c_C(X),\qquad
 (e_*+2)c_B(X^{+8})
 \le100000(e_*+2)+\tfrac12\theta_0c_C(X).
\]
For any prescribed $q$, set
$\Gamma=(q+10)\log H$ and $u_0=(a/2)\log s$.
The stage thresholds can impose
\[
 64(u_0+\Gamma+100000(e_*+2))\le C\theta_0/2,\qquad
 100000C(e_*+2)\le B\theta_0/2.
\]
This is possible because $C$ grows like $(\log s)^2$,
whereas $u_0=O(\log s)$ and $\log H=o(\log s)$.
Consequently, for each term of the original expansion with insertions,
\begin{equation}\label{ch11:opening-higher-weight}
 r^m e^{(e_*+2)c_B(X^{+8})}
 \le e^{-\Gamma}r_0^m e^{\theta_0c_C(X)}.
\end{equation}
This changes weights on the full family of terms of source degree at least two;
it does not compare $r$ directly with $r_0$.
Root-to-site conversion contributes only the factor $\mathcal W(B)=H^{o(1)}$.
At the opening all distinct sites, including those in one
measurement cell, are separated: their distance is a fixed
positive multiple of $H^b$, and $C^2=o(H^b)$.

Choose $U$ so the normalized single-insertion bound at radius $H^f$
is at most $H^{-8}$. Its near Taylor error is at most
$H^{f+bd+c_d-D}$, and its distant-support part is at most
$H^{f+bd+c_d}e^{-\theta_0 C}$. Increase $D$ to make both
at most $H^{-q-2}$. Together with
\eqref{ch11:opening-higher-weight} these give the initial
regular remainder $H^{-q}$ at exponent $e_*$.
For the large-field part perform the exact change of representation at the comparison scale
jointly with these factors with insertions, before summing components.
Every nonzero failed cutoff and every old large-field activity
meets a new large-field site. Its selected site factor has size at
most $e^{-2u}$, with $u=a_*h(s)^2$.
Keep $e^{-u}=G(s)$ aside once per component.
The remaining rooted sum is bounded by
\[
 H^{-q-8}
 +\mathcal W(B)\bigl(O(H^{-8})+O(H^{-q-8})\bigr)
 +H^{o(1)}\sqrt{e^{-H^{1/5}}}
 +(8B+1)^4e^{-u+O(e_*)},
\]
which is at most $H^{-1}$ and satisfies the connected-set
smallness condition after the increases of the starting scales.
Thus its regrouped site sum is at most
$G(s)(e^{H^{-1}}-1)\le G(s)$.
All original and new prescribed large-field sets and compatibility conditions are
retained, and component decomposition is reversible in every
squarefree coefficient. This proves the full opening bound;
no higher coefficient of the transferred density is discarded.

The upper interval has
$m=\lceil(b+2)\log H/\log L\rceil$ steps.
Before choosing $L\ge3$, fix
$\kappa=1+2(b+2)\log(\alpha+2)/\log3$.
Then $(\alpha+2)^m\le H^\kappa$ eventually for every allowed $L$.
Choose $q\ge\kappa+Un+2$ before the Taylor order $D$ and
the later lower bound on the block ratio involving its polynomial transfer norm.
This prevents a circular choice between $q,D$ and $L$.
The lower interval only increases the cell side by a factor
$C^{10}$, so it uses at most
$\lceil10\log C/\log L\rceil$ steps.
Since $\log C=O(\log\log s)=o(\log H)$, its growth factor
$(\alpha+1)^{\ell-w}$ is at most $H^{1/4}$ eventually.
The selected stage plateau must include these two finite windows;
their total length is $O((\log\log s)^2)$ at a fixed stage
and lies inside the short interval admitted by the restart
construction.

Sections~\ref{ch11:weight-restoration} and
\ref{ch11:common-extraction}, together with the restarted
representation, now give the weighted output majorant
\eqref{ch11:weighted-output-envelope}.
To normalize it one must retain the Gaussian power in the
partition function. Axial gauge fixing leaves $3V+1$ free links
on a four-dimensional torus of $V$ sites. Restrict them to a
ball for the normalized Hilbert--Schmidt distance of radius $g/4$ to obtain
\[
 Z\ge e^{\sigma-Vc_1}(c_-g^{\dim\gLie})^{3V+1}.
\]
For the numerator, successive axial plaquette elimination
integrates $3V-3$ distinct links, leaving the four wrapping
links. An explicit elimination order verifies that count.
Represent sites by $0,\ldots,N_{\rm per}-1$ and let $p(x)$
be the largest nonzero coordinate. The tree uses
$(x-e_{p(x)},p(x))$ for $x\ne0$.
Retain the four wrapping links $(-e_i,i)$.
For any other free link $(x,i)$ order it by
\[
 \bigl(t,\mathbf1_{\{i=t\}},\textstyle\sum_jx_j\bigr),
 \qquad t=\max(i,p(x)),
\]
lexicographically. If $i<t$, choose its plaquette in directions
$(i,t)$ based at $x-e_t$. If $i=t$, freeness implies
$x_i=N_{\rm per}-1$; put $w=x+e_i$, take $j=p(w)<i$,
and choose the $(j,i)$ plaquette based at $x-e_j$.
Every other free link in the chosen plaquette has smaller priority.
The chosen plaquettes are distinct, since equality would require
each of two links to have smaller priority than the other.
Integrate in decreasing order. The chosen link occurs once,
as $aUb$ or $aU^{-1}b$, so Haar invariance evaluates its factor.
All other plaquette factors are at most one and may be discarded.
This proves the stated $3V-3$ integrations, including the
wrapping links except for the four explicitly retained ones.
The Haar estimate
$\int_Ge^{-\chi(u)/(g/c)}du\le c_+(g/c)^{\dim\gLie}$
then gives the corresponding upper bound.
It follows that the normalized coefficient with weight $z$
obeys
\[
 z|\mathbb E F|
 \le \exp\bigl(VC_{\rm stage}-4\dim(\gLie)\log g\bigr).
\]
The constant term $\sigma$ cancels. A genuine additive value error
$\delta$ contributes $2V\delta$, one copy in numerator and
one in denominator, and is included in $C_{\rm stage}$.
These estimates use the same nonnegative output density.

At the output $V=q_0$. Its smallest-period site count grows
exponentially in the output logarithmic scale, so
$2\dim(\gLie)\log\beta/q_0\le1$ eventually.
Choose the later lower bound on the starting scale with
$C_{\rm stage}+1\le I_k=\log H/32$.
Applying the preceding estimate to
\eqref{ch11:weighted-output-envelope} gives
\[
 \|\hbox{corrected coefficient }I\|_{\rm root}
 \le e^{I_k}\omega^{-|I|}
 \le e^{-I_k},\qquad \varnothing\ne I\subseteq\{1,2\}.
\]
This is the desired corrected pair estimate on the original
partition denominator. At each fixed deep comparison scale its
simultaneously normalized columns are uniformly bounded in $N,p$.
Undoing the finite connected subtractions and then taking a
sufficiently deep comparison scale establishes
Proposition~\ref{ch11:finite-detector-input}, using the small value and joint product bounds of
Proposition~\ref{ch11:restart-family}. The estimates hold before any subsequence is chosen and their
thresholds are independent of the period.


\subsection{Finite-arity comparison on a common finite set of scales}
\label{ch11:finite-arity-comparison}

The pair argument extends to any fixed arity, with the same
partition normalization and one polynomial localization map for each degree.
We state the result in the form needed for smooth free-field
comparison. All constants may depend on the fixed finite arity,
degree list and compact configuration set.

\begin{proposition}[Common finite-arity Gaussian comparison]
\label{ch11:n-mark-comparison}
Fix $n\ge1$, finitely many real polynomial labels of bounded
degree, and a compact set of point configurations with positive
pairwise distance. Fix finitely many magnification exponents
$b_1,\ldots,b_t$ large enough for the preceding local estimates.
The selected stage requirements can be imposed for this finite
request so that one common localization map $T_{k,N,p,d}$ onto the polynomial label space
per degree and one restarted history serve all these scales.
At scale $\rho_j=H^{b_j}$ use the exact simultaneous lifts
\[
 S_{g/\theta}T_{k,N,p,d}u=W_{\rho_j}v.
\]
They are real and equivariant, and bounded in $N,p$ for each
fixed comparison scale. Set $G_\varnothing=1$ and let $G_I$
be the free moment of the chosen labels at the rescaled
configurations. Define the corrected coefficients, the Wick products of the $F_i$ with respect to the free field, by the recursion
\[
 R_\varnothing=1,\qquad
 R_I=\prod_{i\in I}F_i-
       \sum_{\varnothing\ne J\subseteq I}G_J R_{I\setminus J}.
\]
Then, for every nonempty $I\subseteq\{1,\ldots,n\}$, the
chessboard bound on the corresponding comparison cell is
\begin{equation}\label{ch11:n-mark-root}
 \|R_I\|_{\rm root}\le H^{-1/(16n)}
\end{equation}
at all sufficiently large cutoffs, uniformly in period and
in the stated compact configuration set.
The cell side can be chosen at most $a_kH^{b_j+5}$.
The estimate requires a positive distance between the points, but no
positive lower bound on individual coordinate differences.
\end{proposition}
\begin{proof}
Make all finite requests before choosing the stage orders,
ratios and starting scales. Construct $T$ at the comparison scale $k$ from
the insertion and its original same-insertion covariance.
Neither its near-support Taylor construction nor that covariance
depends on the choice of middle or output scale.
Choose the Taylor order and derivative margins for the
largest of the finitely many requests. Restart once on the
longest required window. Restriction of that history gives
the shorter windows, with unchanged localization map and density at the comparison scale.
There is no new normalization for individual spectator subsets.

Work throughout in the full $n$-source squarefree algebra.
The polynomial $\sum_I G_Iz_I$ has constant coefficient one,
so its inverse exists there. Multiplying the density with insertions
by this inverse gives exactly the displayed recursion.
Equivalently subtract its finite connected logarithm.
This operation commutes with every block integration and leaves
the density without insertions, hence its partition function, unchanged.

Use
\[
 \omega_-=H^{1/(8n)-U},\qquad
 \omega=H^{1/(8n)},\qquad I_k=\frac{\log H}{16n}.
\]
All support and regrouping arguments above already use $n$
only through $n$, $2^n$, and the fixed polynomial degree.
Choose $U\ge1/(8n)$ to control the finite list of single insertions,
then $\kappa,q,D$, and finally the ratios and starting scales in the
order already specified. The initial estimate for terms of source degree at least two
\eqref{ch11:opening-higher-weight} applies to all cardinalities;
the full incoming polynomial is retained.

For the Gaussian calculation under the joint law use degree bound $nD$
and, for example, $\lambda\le1/[32n^2(D+1)]$.
The coefficient tails and the covariance errors on the retained jet block are
$H^{-2}$, so the polynomial Wick estimate makes all subset
moment errors at most $H^{-3/2}$.
The finite logarithm is a polynomial of degree at most $n$
in these subset moments. Its derivative is bounded on a fixed
neighborhood of their free values, since every free subset
moment is bounded on the compact separated configuration set.
Hence the connected coefficient errors are also $O(H^{-3/2})$
and in particular at most $H^{-31/32}$ eventually.
This gives the same corrected principal bound as before.

Each source monomial gains at most
$\omega^{|I|}\le H^{1/8}$ at the larger source weight.
Lemma~\ref{ch11:restore-regular} therefore gives $H^{-1}$;
the corrected principal contribution is at most
$2^nR_eH^{-27/32}$; and the regrouping of activities constrained to cover their prescribed large-field sets has bound
$4H^{Un+2^n}\eta_-$.
The upper and lower recurrences then give incidence and site
bounds at most one at the output. Their scalar majorant contributes
at most $1+n$ per site. The partition
estimate absorbs this finite factor by imposing
$C_{\rm stage}+1\le I_k$.
Thus the corrected coefficient has chessboard seminorm at most
$e^{I_k}\omega^{-|I|}\le e^{-I_k}$, proving
\eqref{ch11:n-mark-root}.

For the geometric bound, the scale ratio of the upper interval is at most
$L H^{b_j+2}$ and the lower ratio at most $L C^{10}$.
The lower bounds on the starting scales give $L^2C^{10}\le H^3$, so the total
cell side in continuum units is at most $a_kH^{b_j+5}$.
All covariance estimates depend only on the full displacement
between distinct points. On the compact configuration set its minimum
length is positive, also after rounding at sufficiently large
$H$. The opening-support estimate uses the same positive
maximum-coordinate distance. Neither argument divides by a
single coordinate difference. This proves the asserted
uniformity, including configurations arbitrarily close to
coordinate coincidences.
\end{proof}

Different scales in this proposition use their own normalized
lifts, but the coefficient map is identical. Consequently, if
$b_0>b_j$, the lift normalized at $b_0$ is exactly the linear
combination of the degree-$d$ lifts at $b_j$ obtained by
applying $W_{H^{b_0-b_j}}$ to its free label.
On a fixed finite-dimensional degree space that contributes a factor at most
$C H^{(b_0-b_j)d}$ in coefficient norm.
The chessboard estimate for a single insertion and its free expectation therefore
give the same factor for its chessboard seminorm.
This is the comparison needed for the size and displacement
bounds at three scales; it does not assert that one fixed
unnormalized label has a small corrected error at all three.


We can now state the finite-cutoff comparison given by these estimates.

\begin{proposition}[Comparison of the separating correlations at finite cutoff]
\label{ch11:finite-detector-input}
On the common trajectory with the stage orders, lower bounds on the block ratio and
choices of sufficiently large starting scales specified above, for each $e$ there is an injective free measurement matrix
$G:\CC^r\to\CC^s$ and the following finite-cutoff approximations.
For arbitrarily small $\eta>0$ one can choose a fixed comparison
scale $k$, bounded real local polynomial coefficient families
$u_{1,i},\ldots,u_{r,i}$ with $\|u_{j,i}\|_1\le B_k$, and finitely
many bounded families of normalized local spectator insertions, all on the selected
history, such that their point-correlation matrix $Q_i^0$
satisfies
\[
 \|Q_i^0-G\|_{\max}\le\eta
\]
eventually for every period. The correlation configurations and small
neighborhoods of them have a positive common coordinate distance
and lie inside the smallest period in continuum units. All these bounded
normalized local insertion families have uniform chessboard bounds at every
sufficiently fine fixed cell depth.
\end{proposition}
\begin{proof}
Choose the free basis and separating configurations from
Chapter~\ref{ch03:maxwell}. Construct their common
localization map and simultaneous lifts by
Section~\ref{ch11:common-extraction}, and use the restarted family
of Proposition~\ref{ch11:restart-family}.
The opening estimate, stronger support bounds, larger source weight and output normalization in
Section~\ref{ch11:corrected-pair-comparison} give corrected
subset chessboard-seminorm error $e^{-I_k}$, uniformly in cutoff and period.
Undoing the connected subtractions gives point-matrix error at
most $C_ke^{-I_k}$, where $C_k$ is a polynomial in the finitely
many free subset moments. These moments remain bounded as the
rounded, rescaled configurations approach their fixed separated
limit configurations. Thus choose one sufficiently deep $k$ to obtain the
requested $\eta$. The exact lift equations give a fixed bound
$B_k$ for all cutoffs and periods at that scale.
The product-moment estimate gives the chessboard bounds for normalized local insertions.
No convergence or choice of a limiting equation subspace is used.
\end{proof}

This is a comparison before any cutoff limit. Its common localization
map, auxiliary density and normalized lifts are the ones
constructed above.

Recall the exact cutoff-dependent equation kernel
$K_i=\mathsf Z_{N,p,e}P_{N,p,e}S_{t_N}R_e\mathcal I_e$.
Its bounded projection $\Pi_i$ is distinct from the near-identity
section comparison $P_{N,p,e}$. The latter is invertible; the
former has kernel $K_i$.

\begin{proposition}[A separating correlation family on the entire cutoff-dependent complement]
\label{ch11:uniform-detector}
Using Proposition~\ref{ch11:finite-detector-input}, for every fixed
finite difference word $l$ there are averaged separating correlation functionals $L_i$ and
a constant $Z$, independent of $N,p$, such that
\[
 \|\Pi_i v\|_1\le Z\|L_i\Pi_i v\|_\infty
 \quad(v\in V_e)
\]
eventually for every period. Their values on a finite difference
$D_lF$ of a normalized local insertion with bounded chessboard seminorm are uniformly bounded.
Their values on any trial whose chessboard seminorm tends to zero tend to
zero uniformly in period.
\end{proposition}
\begin{proof}
Choose $K$ with $\|c\|_\infty\le K\|Gc\|_\infty$ and put
$\delta=(2Kr)^{-1}$. The case $r=0$ is immediate.
Use the assumed comparison with entrywise error at most $\delta/4$.
The modulus of continuity in the insertion points for normalized local insertions from
Section~\ref{ch11:positional} holds on the separated neighborhoods
with one constant and threshold for this finite list.

Choose $b>0$ so small that moving the distinguished point by at most
$2|l|b$ stays in those neighborhoods and changes every normalized local insertion
measurement by at most $\delta/4$. Put $h_N=\lceil b/a_N\rceil$.
Average over the shifts $0,\ldots,h_N-1$ once in each coordinate
listed in $l$, always in the distinguished slot. These normalized
averages define $L_i$. For the empty word take no averages.

The vectors $(I-\Pi_i)u_{j,i}$ are bounded, hence admissible.
Lemma~\ref{ch11:uniform-sweep} and the chessboard inequality with one distinguished insertion
show that the values of the averaged functionals on them tend to zero.
Thus the matrix
\[
 Q_i=(L_i\Pi_i u_{j,i})_{j=1}^r
\]
is within $\delta$ of $G$ eventually. For any $c\in\CC^r$,
\[
 \|c\|_\infty
 \le K\|Q_ic\|_\infty+Kr\delta\|c\|_\infty,
 \qquad
 \|c\|_\infty\le2K\|Q_ic\|_\infty.
\]
The $r$ vectors $\Pi_i u_{j,i}$ are therefore independent.
They span $\operatorname{ran}\Pi_i$, since its dimension is
exactly $r$. Writing $\Pi_i v=\sum_jc_j\Pi_i u_{j,i}$ gives
\[
 \|\Pi_i v\|_1\le ArB_k\|c\|_\infty
       \le2ArB_kK\|L_i\Pi_i v\|_\infty.
\]
Take $Z\ge\max(1,2ArB_kK)$.

For the assertions on the values of the averaged functionals, inserting a distinguished trial of chessboard seminorm $u_i$ with the bounded spectator fields gives a bound $Cu_i$
uniformly throughout the neighborhoods. Averaging preserves it.
For a difference word, differences commute with these averages,
and each coordinate telescopes exactly:
\[
 \frac1{h_N}\sum_{j=0}^{h_N-1}
 \frac{H(x+(j+1)e_\alpha)-H(x+je_\alpha)}{a_N}
 =\frac{H(x+h_Ne_\alpha)-H(x)}{h_Na_N}.
\]
Iteration bounds the averaged difference by
$(2/b)^{|l|}$ times the undifferentiated correlation bound.
All endpoints stay in the chosen neighborhoods.
This proves the bound on these values without assuming any pointwise
bound on the differentiated field.
\end{proof}

\begin{theorem}[Uniform bounds for finite differences]
\label{ch11:all-period-differences}
For the common trajectory with the trial estimates of
Section~\ref{ch11:trial-descent} and the stage requirements already
specified for their uniformity, every fixed finite list of
fields, every finite difference word applied separately to its
slots, and every $r_0>0$, their finite-volume correlation is bounded
by one constant for all sufficiently large $N$, every period $p$,
and every tuple with pairwise periodic maximum-coordinate distance
at least $r_0$ in continuum units.
\end{theorem}
\begin{proof}
Start with a bounded local polynomial coefficient and a difference word $l$.
The exact physical difference identity and the translated local
section bound make this a trial at degree $e=d+|l|$ satisfying the hypotheses of \eqref{ch11:trial-error}.
The factor $a_N^{|l|}$ in the normalization at its filtration index cancels all inverse mesh
powers. Its link support uses at most $2^{|l|}$ fixed translates;
its derivative bound grows by at most
$(1+2^{\varpi_Nq_N})^{|l|}$, a subpower quantity under the stage
requirement on the starting scale. The backward comparison gives an admissible coefficient
$x_i$ and an error $E_i=O(x_i)-D_lF$ with chessboard seminorm tending to zero.

The correlation pairings just proved give $\|L_ix_i\|\le C$.
Apply Lemma~\ref{ch11:absorption} with
$y_i=\Pi_i x_i$ and
$z_i=(x_i-y_i)/(1+\|y_i\|_1)$.
The latter family is admissible: its first term is a bounded
scalar multiple of $x_i$, and its second is norm bounded.
It lies in the exact kernel of $\Pi_i$, so the values of the correlation functionals on it tend to
zero by Lemma~\ref{ch11:uniform-sweep}.
Consequently $\|y_i\|_1\le2ZC+1$.
The chessboard bound for normalized local insertions applies to $O(y_i)$,
and the unnormalized admissible kernel family $x_i-y_i$
has vanishing chessboard seminorm. The identity
\[
 D_lF=O(y_i)+O(x_i-y_i)-E_i
\]
gives a uniform chessboard bound for the difference trial.
Repeat for the finite list of orientations and slots.

For clarity, two cell scales are needed to place a separated tuple in distinct cells.
Choose a fixed coarse cell side $a_\ell<r_0$.
Choose a finer fixed depth $k\ge\ell$, beyond all the finitely many
scale thresholds for the chessboard bounds, with $6na_k<a_\ell$.
At mesh $N$ the corresponding integer cell sides $t,T$ satisfy
$a_Nt=a_k$ and $a_NT=a_\ell$.
In each coordinate one point forbids at most $6t$ translations
of the coarse tiling if a margin $3t$ is required.
The $n$ points therefore leave a permissible translation, since
$6nt<T$. Distinct points occupy distinct coarse cells: sharing
a cell would make their periodic maximum-coordinate distance
at most $a_\ell<r_0$.
The chessboard-seminorm comparison between aligned fine and coarse
tilings transfers the depth-$k$ bounds to these interior cells;
it follows from the repeated chessboard inequality (maximizing
the position of a repeated copy in its containing cell).
Translation invariance alone is not this comparison.
Placing the finitely many trials in distinct cells now bounds the correlation by
the product of their fixed chessboard bounds.
All cell sides, margins and lower bounds on $N$ were chosen independently
of $p$. The periodic difference identity also holds at wrapped
sites, completing the proof.
\end{proof}

